% Source: preamble
\documentclass[11pt,letterpaper]{article}
\usepackage[margin=1.05in]{geometry}
\usepackage[T1]{fontenc}
\usepackage{lmodern,microtype}
\usepackage{amsmath,amssymb,amsthm,mathtools}
\usepackage{xcolor}
\definecolor{paperlink}{RGB}{0,0,114}
\usepackage[colorlinks=true,hypertexnames=false,linktoc=all,linkcolor=paperlink,citecolor=paperlink,urlcolor=paperlink]{hyperref}
\usepackage{fancyhdr}
\pagestyle{fancy}
\fancyhf{}
\fancyhead[L]{\small Gaussian approximation over rectangles}
\fancyhead[R]{\small\thepage}
\setlength{\headheight}{14pt}
\hypersetup{pdftitle={High-Dimensional Gaussian Approximations over Rectangles without Covariance Restrictions},pdfsubject={Bounded independent summands; rates optimal up to logarithmic factors},pdfauthor={Human Collaborators: CCKK (Victor Chernozhukov, Denis Chetverikov, Kengo Kato, Yuta Koike); AI Collaborators: GPT-6 Sol and GPT-6 Astra}}
\setlength{\parindent}{1.25em}
\setlength{\parskip}{0pt}
\emergencystretch=2em
\allowdisplaybreaks[2]
\raggedbottom
\clubpenalty=10000
\widowpenalty=10000
\newtheorem{theorem}{Theorem}
\newtheorem{proposition}{Proposition}
\newtheorem{lemma}{Lemma}
\newtheorem{corollary}{Corollary}
\theoremstyle{remark}
\newtheorem{remark}{Remark}
\newcommand{\E}{\mathbb E}
\newcommand{\Pp}{\mathbb P}
\newcommand{\R}{\mathbb R}
\newcommand{\cR}{\mathcal R}
\newcommand{\cF}{\mathcal F}
\newcommand{\cP}{\mathcal P}
\newcommand{\cD}{\mathcal D}
\newcommand{\Cov}{\operatorname{Cov}}
\newcommand{\Var}{\operatorname{Var}}
\newcommand{\tr}{\operatorname{tr}}
\newcommand{\op}{\mathrm{op}}
\newcommand{\norm}[1]{\left\lVert#1\right\rVert}
\newcommand{\rh}{\rho_{\mathcal R}}
\newcommand{\ind}{\mathbf 1}
\newcommand{\dd}{\,d}
\newcommand{\Hh}{H_2}
\newcommand{\Ga}{\gamma_r}
\newcommand{\ph}{\varphi_r}
\newcommand{\surf}{\mathcal J}

\newcommand{\logp}{\log_+}


% Source: front
\begin{document}
\begin{center}
{\Large High-Dimensional Gaussian Approximations\\[0.15em]
over Rectangles without Covariance Restrictions}\\[0.85em]
{\normalsize\textbf{Human Collaborators: CCKK}}\\[0.25em]
{\normalsize Victor Chernozhukov, Denis Chetverikov, Kengo Kato, Yuta Koike}\\[0.65em]
{\normalsize\textbf{AI Collaborators:} GPT-6 Sol and GPT-6 Astra}\\[0.45em]
{\normalsize 21 October 2026}
\end{center}
\thispagestyle{plain}

\begin{abstract}
We study Gaussian approximation over rectangles for sums of independent
bounded random vectors, without a local nondegeneracy assumption on the
covariance matrix. After marginal standardization, consider normalized
observations with a fixed envelope \(B\), and put \(L=\log(ep)\). We prove
\[
 \rho_{\mathcal R}(W_n,Z_n)
 \le \frac{C_B L^{3/2}}{\sqrt n}
       \sqrt{1+\log_+\sqrt{n/L}},
 \qquad W_n=n^{-1/2}\sum_iX_i,
\]
where \(Z_n\) has the same covariance as \(W_n\). If the correlation
gap is bounded below by \(c_0>0\), this improves to
\[
 \frac{C_{B,c_0}}{\sqrt n}
 \left[L^{3/2}+L\sqrt{1+\log_+\frac{\sqrt n}{L}}\right].
\]
Thus \(p\asymp n^a\), with fixed \(a>0\), gives rates \(n^{-1/2}(\log n)^2\) and
\(n^{-1/2}(\log n)^{3/2}\), respectively. The unrestricted rate is
optimal up to logarithmic factors; the power \(n^{-1/2}\) is sharp even
in one dimension. If the full aggregate third-moment tensor vanishes,
the \(L^{3/2}\) term can be omitted from the fixed-gap bound. At a fixed
positive gap, the resulting dimension thresholds
\(\log p=o(n^{1/3})\) and \(\log p=o(\sqrt n)\) are sharp for the
respective bounded-observation classes. The constants are independent
of covariance rank and eigenvalues. The necessity of the extra
unrestricted square-root logarithm remains open. A clipping argument
extends the general bounds to sub-exponential coordinates with an
additional \(\log(en)\) factor.
\end{abstract}

\section{Main result}\label{sec:results}
Let \(\mathcal R\) be the axis-aligned rectangles in \(\R^p\), and write
\[
 \rh(U,V)=\sup_{A\in\mathcal R}
       |\Pp(U\in A)-\Pp(V\in A)|.
\]
This distance controls maxima and simultaneous coordinatewise confidence
regions. Throughout, let
\begin{equation}
 \begin{gathered}
 W=\sum_{i=1}^n\xi_i,\qquad
 \E\xi_i=0,\qquad \norm{\xi_i}_\infty\le\delta\quad\text{a.s.},\\
 \Sigma=\Cov(W),\qquad\Sigma_{jj}=1\ (j\le p),
 \qquad Z\sim N(0,\Sigma).
 \end{gathered}
 \label{eq:setup}
\end{equation}
where the rows are independent. Coordinates within a row may be
dependent, and the rows need not have the same law. Set
\[
 L=\log(ep),\qquad
 D_L=1+\log_+\frac1{\delta\sqrt L},\qquad
 \log_+x=\max\{0,\log x\}.
\]
Constants denoted by \(C\) are absolute unless a subscript records
their dependence.

\begin{theorem}[Unrestricted covariance]\label{thm:unrestricted}
Under~\eqref{eq:setup},
\begin{equation}
 \rh(W,Z)\le C\min\{1,\delta L^{3/2}\sqrt{D_L}\}.
 \label{eq:unrestricted}
\end{equation}
\end{theorem}

The covariance rank is unrestricted, and no lower bound is imposed on
positive eigenvalues or local determinants beyond the marginal
normalization. In particular, singular covariance and arbitrarily strong
correlations are permitted.

\begin{corollary}[Normalized observations]\label{cor:observations}
Let \(X_1,\ldots,X_n\) be independent and centered. Suppose
\(\norm{X_i}_\infty\le B\) almost surely and
\(n^{-1}\sum_i\E X_{ij}^2=1\) for every \(j\). Then, with
\(W_n=n^{-1/2}\sum_iX_i\) and \(Z_n\sim N(0,\Cov(W_n))\),
\begin{equation}
 \rh(W_n,Z_n)\le C\min\!\left\{1,
 \frac{B L^{3/2}}{\sqrt n}
 \sqrt{1+\log_+\frac{\sqrt n}{B\sqrt L}}\right\}.
 \label{eq:observations-unrestricted}
\end{equation}
\end{corollary}

This is Theorem~\ref{thm:unrestricted} with
\(\xi_i=X_i/\sqrt n\). Unequal marginal variances are handled by
positive diagonal scaling: if
\(\sigma_j^2=n^{-1}\sum_i\E X_{ij}^2\),
\(\underline\sigma=\min_j\sigma_j>0\), and
\(\max_i\norm{X_i}_\infty\le K_n\), use the envelope
\(\delta=K_n/(\underline\sigma\sqrt n)\). Rectangles and rectangle
distance are preserved by this scaling.

\subsection{Correlation gaps and third-moment cancellation}
The main result has a stronger form that records an optional improvement
from the covariance. Define
\[
 c=1-\max_{j\ne k}|\Sigma_{jk}|,
 \qquad c=1\ \text{when }p=1,
\]
and put
\begin{equation}
 \kappa=\begin{cases}\min\{L,c^{-1}\},&c>0,\\L,&c=0,\end{cases}
 \qquad E_\kappa=1+\log_+\frac{\sqrt\kappa}{\delta L}.
 \label{eq:parameters}
\end{equation}
A positive gap excludes equal or opposite coordinates; it still permits
singular covariance matrices.

\begin{theorem}[Gap-sensitive bound]\label{thm:optimized}
Under~\eqref{eq:setup},
\begin{equation}
 \rh(W,Z)\le C\min\!\left\{1,
 \delta[L^{3/2}+L\sqrt{\kappa E_\kappa}]\right\}.
 \tag{R*}\label{eq:optimized}
\end{equation}
\end{theorem}

Since \(\kappa\le L\) and \(E_\kappa\le D_L\), this implies
Theorem~\ref{thm:unrestricted}. The \(\delta L^{3/2}\) term comes
from the cubic part of the comparison. It disappears under cancellation
of the full third-moment tensor.

\begin{theorem}[Cancellation of third moments]\label{thm:zero-skew}
Under~\eqref{eq:setup}, suppose
\begin{equation}
 \sum_i\E\xi_i^{\otimes3}=0,
 \quad\text{equivalently,}\quad
 \sum_i\E[\xi_{ij}\xi_{ik}\xi_{i\ell}]=0
 \quad(j,k,\ell\le p).
 \label{eq:zero-third}
\end{equation}
Then
\begin{equation}
 \rh(W,Z)\le C\min\{1,\delta L\sqrt{\kappa E_\kappa}\}.
 \tag{ZS}\label{eq:zero-skew}
\end{equation}
\end{theorem}

Central symmetry of each row is sufficient. Vanishing marginal third
moments alone is insufficient, because the condition includes mixed
third moments. It allows cancellation between different rows and is
preserved by deterministic linear maps.

For a fixed standardized envelope and \(c\ge c_0>0\), the sample-size
bounds become
\begin{align}
 \rh(W_n,Z_n)&\le C_{B,c_0}\min\!\left\{1,
 \frac{L^{3/2}}{\sqrt n}
 +\frac L{\sqrt n}\sqrt{1+\log_+\frac{\sqrt n}{L}}\right\},
 \label{eq:observations-fixed}\\
 \rh(W_n,Z_n)&\le C_{B,c_0}\min\!\left\{1,
 \frac L{\sqrt n}\sqrt{1+\log_+\frac{\sqrt n}{L}}\right\}
 \quad\text{under~\eqref{eq:zero-third}}.
 \label{eq:zero-fixed}
\end{align}
These bounds give uniform convergence at a fixed positive gap when
\(\log p=o(n^{1/3})\) in general and \(\log p=o(\sqrt n)\)
under cancellation. Theorem~\ref{opt:consistency} shows that these
conditions are sharp over fixed-envelope classes, with envelope at least
two for the first assertion and at least one for the second. The latter is a sharp dimension threshold, not a sharp error
rate below that threshold.

\paragraph{Lean certification.}
The bounded upper bounds in Theorems~\ref{thm:unrestricted}--\ref{thm:zero-skew}
were first certified on 20 October 2026 using Lean 4.19.0 and
mathlib v4.19.0. The certified Git revision is
\href{https://github.com/VC2015/HDCLTs/commit/3804008bc7ee4bbde244f8f46e9d4e8a58103185}{\texttt{3804008bc7ee}};
the combined Lean theorem is
\nolinkurl{RectangleCLT.Companion.companionCLTExists_common}.
The transitive axiom audit reports only Lean's standard axioms; the final
theorems have no unproved analytic estimates as premises.\footnote{Full Git SHA:
\nolinkurl{3804008bc7ee4bbde244f8f46e9d4e8a58103185}. The
\href{https://github.com/VC2015/HDCLTs/blob/3804008bc7ee4bbde244f8f46e9d4e8a58103185/rectangle-clt-singular/VERIFICATION.json}{verification record}
identifies the toolchain, source hashes, and successful
build and audit commands.}
The certification covers these three bounded-observation upper bounds;
the sub-exponential extension and the supplementary lower bounds and
consistency thresholds are outside its scope.

\subsection{Sub-exponential coordinates}\label{sec:subexponential}
The bounded theorem extends by clipping the coordinates and controlling
the discarded sum. We use
\[
 \|V\|_{\psi_1}=\inf\{t>0:\E e^{|V|/t}\le2\}.
\]
The cutoff depends on \(\log n\), rather than \(\log(np)\).

\begin{corollary}[Sub-exponential observations]\label{cor:subexponential}
Let \(X_1,\ldots,X_n\in\R^p\) be independent and centered, with
\(n^{-1}\sum_i\E X_{ij}^2=1\) and
\(\max_{i,j}\|X_{ij}\|_{\psi_1}\le K\), where \(K\ge1\).
For \(W_n=n^{-1/2}\sum_iX_i\) and a centered Gaussian \(Z_n\)
with the same covariance, put
\[
 \ell_n=\log(en),\qquad b_n=\frac{K\ell_n}{\sqrt n},\qquad
 D_n=1+\log_+\frac1{b_n\sqrt L}.
\]
Then
\begin{equation}
 \rh(W_n,Z_n)\le C\min\{1,b_nL^{3/2}\sqrt{D_n}\}.
 \label{eq:subexponential}
\end{equation}
If the correlation gap of \(W_n\) is at least a fixed \(c_0>0\), then
\begin{equation}
 \rh(W_n,Z_n)\le C_{c_0}\min\!\left\{1,
 b_n\left[L^{3/2}+L\sqrt{1+\log_+\frac1{b_nL}}\right]\right\}.
 \label{eq:subexponential-gap}
\end{equation}
No local determinant or positive-eigenvalue lower bound is needed for
\eqref{eq:subexponential}. Unequal marginal variances are handled by
imposing the tail condition on \(X_{ij}/\sigma_j\), as in the
standardization following Corollary~\ref{cor:observations}.
\end{corollary}

For fixed \(K\) and \(p\asymp n^a\), \(a>0\), these rates are
\(n^{-1/2}(\log n)^3\) and
\(n^{-1/2}(\log n)^{5/2}\), respectively. Thus the sample-size
power remains optimal up to logarithmic factors. In particular,
\eqref{eq:subexponential} guarantees convergence for fixed \(K\) when
\(\log p=o(n^{1/3}/\log(en))\). These sufficient dimension ranges
are not claimed sharp for the sub-exponential class.
Appendix~\ref{app:subexponential} controls the discarded sum directly,
which allows the cutoff to depend only on \(\log n\), and restores
covariance through a joint Gaussian coupling.

\subsection{Relation to recent work}\label{sec:literature}
High-dimensional Gaussian approximation for maxima and rectangles was
developed by Chernozhukov, Chetverikov, and Kato~\cite{CCK2013,CCK2017}.
Without covariance restrictions beyond positive marginal variances,
Chernozhukov, Chetverikov, Kato, and Koike~\cite[Theorem~2.1]{CCKK2022}
obtained an \(n^{-1/4}\) rate, up to logarithms, under coordinatewise
sub-exponential tails and an averaged fourth-moment condition.
Gao and Zhang~\cite[Theorem~1 and Section~5]{GaoZhang} recently improved
the sample-size power to \(n^{-1/3}\). Under the tail assumption of
Corollary~\ref{cor:subexponential}, their rectangle bound, in our
unit-variance notation, is
\begin{equation*}
 C\min\left\{1,
 \left(\frac{K^2}{n}\right)^{1/3}\{\log(2np)\}^{7/3}
 +\frac K{\sqrt n}\{\log(2np)\}^{5/2}\right\}.
\end{equation*}
No additional fourth-moment condition is needed for this bound.
For fixed \(K\) and \(p\asymp n^a\), with fixed \(a>0\), it has
order \(n^{-1/3}(\log n)^{7/3}\). In this dimension regime,
Corollaries~\ref{cor:observations} and~\ref{cor:subexponential}, with
fixed \(B\) or \(K\), give the sharp sample-size power
\(n^{-1/2}\), with logarithmic factors \((\log n)^2\) for bounded
observations and \((\log n)^3\) for sub-exponential coordinates.
These comparisons concern fixed envelopes or tail scales; the full
bounds retain their different dependence on these parameters.

A second line of work obtains rates close to \(n^{-1/2}\) under
covariance nondegeneracy. For bounded sums with positive definite
\(\Sigma\), Aronow and Lopatto~\cite[Theorem~1.1]{AL} remove the
\(\log(en)\) factor from the bound of Chernozhukov, Chetverikov,
and Koike~\cite[Corollary~2.1]{CCK}, and prove
\begin{equation*}
 \rh(W_n,Z_n)\le C\min\left\{1,
 \frac{BL^{3/2}}{\lambda_{\min}(\Sigma)\sqrt n}\right\}.
\end{equation*}
A fixed positive minimum eigenvalue therefore gives the sharp
\(BL^{3/2}/\sqrt n\) scale in the regimes of matching lower bounds.
Our unrestricted bound remains finite as eigenvalues vanish, at a
square-root logarithmic cost relative to this scale.

Full and local nondegeneracy also lead to different results.
Fang, Koike, Liu, and Zhao~\cite[Theorem~1.1]{FKLZ} obtain nearly
\(n^{-1/2}\) rates under moment or tail assumptions and quantitative
two- and three-coordinate covariance conditions, allowing the full
covariance matrix to be singular. The companion paper~\cite{Local}
assumes uniformly positive eigenvalues for principal submatrices of
order at most four. For bounded observations it gives
\(C_aBL^{3/2}/\sqrt n\), and \(C_aBL/\sqrt n\) under aggregate
third-moment cancellation. The present results also cover singular
three- and four-coordinate submatrices. Thus the two papers give
complementary bounds, with the remaining unrestricted logarithmic loss
isolated in \(\sqrt{D_L}\).

There is a close connection between the proofs as well. Gao and Zhang
use support reduction, Bernoulli interpolations, and rank-independent
Gaussian matrix-score and polytope surface estimates
\cite[Proposition~4 and Lemma~12]{GaoZhang}. Our argument combines
the rectangle concentration input of Aronow and Lopatto with surface
estimates averaged over a concentrated base. An exact expansion of
translated selector weights then controls the fourth-order remainder
along one covariance-preserving path, while retaining aggregate
third-moment cancellation. Appendix~\ref{appx-sec:support} gives these
estimates; Appendix~\ref{sec:main-proof} explains the moving cutoff
that produces the stated logarithmic factor.

The next section explains the two analytic inputs: smoothing and a
fourth-order comparison. Appendices~\ref{appx-sec:smoothing}--\ref{app:subexponential}
contain all proofs of the upper bounds. The Supplement collects the
lower-bound calculations, exact coordinate reductions, and further
investigations.


% Source: proof_interface
\section{Smoothing and covariance-preserving comparison}
\label{sec:smoothing}\label{sec:support}

We first replace a rectangle indicator by a rational selector, then
compare its expectations along a path that preserves covariance.
The analytic estimates are stated here; their proofs are in
Appendices~\ref{appx-sec:smoothing}--\ref{appx-app:concentration}.
Appendix~\ref{sec:main-proof} integrates and optimizes the comparison.

For signed coordinate faces
\(A=\{x:\nu_\alpha\cdot x\le a_\alpha,\ \alpha\le m\}\), define
\begin{equation}
 F_{A,h}(x)=\left(1+\sum_{\alpha=1}^m
          e^{(\nu_\alpha\cdot x-a_\alpha)/h}\right)^{-1},
 \qquad M_h=\sup_A|\E F_{A,h}(W)-\E F_{A,h}(Z)|.
 \label{eq:selector}
\end{equation}
The face normals are distinct members of \(\{\pm e_1,\ldots,\pm e_p\}\),
so \(m\le2p\). The same estimates hold for inconsistent threshold
systems, which can arise when the faces are shifted independently.

The following concentration input controls both the smoothing error
and the derivatives near the boundary. For
\(A=\prod_j[\ell_j,u_j]\), write
\(A^a=\prod_j[\ell_j-a,u_j+a]\), with reversed intervals empty.
\begin{lemma}[Rectangle concentration]\label{lem:concentration}
Let \(S=\sum_i\zeta_i\) have independent centered summands with
\(\norm{\zeta_i}_\infty\le\eta\) and
\(\Var(S_j)=v>0\) for every coordinate. For an independent shift
\(R\), a rectangle \(A\), and \(a>0\),
\[
 \Pp(S+R\in A^a\setminus A^{-a})
 \le C\min\{1,(a+\eta)\sqrt{\log(ep)/v}\}.
\]
For a centered Gaussian base with equal coordinate variances \(v\),
the same bound holds with \(\eta=0\), including singular covariance.
\end{lemma}
Its proof, following~\cite[Lemma~2.1]{AL}, is given in
Appendix~\ref{appx-app:concentration}. Applied to signed face
coordinates, it also handles empty inverse images of threshold systems.

\begin{lemma}[Smoothing error]\label{lem:endpoint}
Under~\eqref{eq:setup}, for \(h\ge\delta\),
\begin{equation}
 \E|F_{A,h}(W)-\ind_A(W)|+
 \E|F_{A,h}(Z)-\ind_A(Z)|\le Ch\sqrt L\,\kappa.
 \label{eq:endpoint-error}
\end{equation}
Consequently,
\begin{equation}
 \rh(W,Z)\le M_h+Ch\sqrt L\,\kappa.
 \label{eq:endpoint-min}
\end{equation}
\end{lemma}
The unrestricted error is \(ChL^{3/2}\). Indeed, independent unit
exponentials give the exact mixture
\[
 F_{A,h}(x)=\Pp_E\{\nu_\alpha\cdot x\le
 a_\alpha+h\log(E_\alpha/E_0)\text{ for every }\alpha\}.
\]
The expected largest threshold displacement is \(O(hL)\), and
Lemma~\ref{lem:concentration} supplies the factor \(\sqrt L\).
When \(c>0\), a Hessian estimate for
\(h\log(1+\sum_\alpha e^{(\nu_\alpha\cdot x-a_\alpha)/h})\)
gives the alternative error \(Ch\sqrt L/c\). Choosing the smaller
estimate explains \(\kappa\). Lemmas~\ref{appx-lem:endpoint-free}
and~\ref{appx-lem:endpoint-gap} prove the two endpoint bounds.

For the smooth comparison, let the Gaussian \(Z\) be independent of
the rows. For \(\lambda\in[0,1]\), set
\begin{equation}
 a=1+\lambda,\qquad q=\frac{1-\lambda}{(1+\lambda)^2},\qquad
 U_\lambda=a\sum_i B_{i,\lambda}\xi_i+\sqrt\lambda Z,
 \qquad B_{i,\lambda}\sim\operatorname{Bernoulli}(q),
 \label{eq:path}
\end{equation}
with independent Bernoulli variables, rows, and Gaussian vector.
Since \(qa^2=1-\lambda\), this path has covariance \(\Sigma\),
with \(U_0\stackrel d=W\) and \(U_1\stackrel d=Z\).
For \(P_A(\lambda)=\E F_{A,h}(U_\lambda)\), a fourth-order expansion gives
\begin{equation}
 P_A'(\lambda)=-\frac\lambda3\,
 \E D^3F_{A,h}(U_\lambda):\sum_i\E\xi_i^{\otimes3}
 +R_{4,A,h}(\lambda).
 \label{eq:path-summary}
\end{equation}
The colon denotes contraction of all three indices. The cubic tensors
are evaluated against the same derivative tensor at \(U_\lambda\).
Under \eqref{eq:zero-third}, they therefore cancel before an absolute
value is taken, and
\begin{equation}
 P_A'(\lambda)=R_{4,A,h}(\lambda).
 \label{zs:identity}
\end{equation}
Moving derivatives from deleted-row bases back to \(U_\lambda\)
creates fourth-order terms, which remain in \(R_{4,A,h}(\lambda)\).

\begin{proposition}[Two time ranges]\label{an:prop:derivative}
There are sufficiently large absolute constants \(A,d_0\) such that,
for \(h\ge A\delta\), uniformly over the thresholds,
\begin{align}
 |P_A'(\lambda)|&\le C\sqrt L\left(
 \frac{\lambda\delta}{h^2}+\frac{\delta^2}{h^3}\right),
 &&0<\lambda<1,\label{an:early}\\
 |P_A'(\lambda)|&\le CL^{3/2}\left(
 \delta+\frac{\delta^2}{\lambda h}\right),
 &&d_0\delta^2L\le\lambda<1.\label{an:late}
\end{align}
The fourth-order remainder separately satisfies
\begin{align}
 |R_{4,A,h}(\lambda)|&\le C\delta^2\sqrt L/h^3,
 &&0<\lambda<1,\label{eq:remainder-early}\\
 |R_{4,A,h}(\lambda)|&\le C\delta^2L^{3/2}/(\lambda h),
 &&d_0\delta^2L\le\lambda<1.\label{an:remainder}
\end{align}
\end{proposition}
The path is continuous in law at both endpoints, and the selector is
bounded and continuous. Together with the integrable early derivative
bound, this justifies integration up to \(0\) and \(1\).
The early estimates use boundary concentration. The later estimates
use Gaussian integration by parts on the covariance support, averaged
surface bounds, and restoration of the actual translated row bases.
Working on the support avoids any lower bound on a positive covariance
eigenvalue. The exact remainder kernels and these analytic verifications
are supplied in Appendix~\ref{appx-sec:support}.


% Source: rate_explanation
\section{From the comparison to the rate}
The early and late fourth-order bounds agree at the scale
\(\tau_h=Lh^2\). If \(\tau_h<1\), split the path at this time.
Choosing \(A^2\ge d_0\) ensures that \(h\ge A\delta\) places
the split in the range of the late estimate. Integrating gives
\[
 M_h\le CL^{3/2}\left[
 \delta+\frac{\delta^2}{h}\log_+\frac1{Lh^2}\right],
 \qquad h\ge A\delta.
\]
When \(\tau_h\ge1\), the early estimate suffices for the whole path.
Under aggregate third-moment cancellation, the cubic term is absent,
and the comparison bound becomes
\[
 M_h\le\frac{C\delta^2L^{3/2}}h
 \left[1+\log_+\frac1{Lh^2}\right].
\]
The smoothing error is \(Ch\sqrt L\,\kappa\). Balancing it against
the comparison remainder leads to
\[
 h=A\delta\sqrt{LE_\kappa/\kappa}.
\]
Since \(\kappa\le L\) and \(E_\kappa\ge1\), this scale is at
least \(A\delta\). Also
\(1+\log_+(1/(Lh^2))\le CE_\kappa\), so both the smoothing
error and the fourth-order comparison error are at most
\(C\delta L\sqrt{\kappa E_\kappa}\).
This yields Theorems~\ref{thm:optimized} and~\ref{thm:zero-skew};
Theorem~\ref{thm:unrestricted} follows from \(\kappa\le L\) and
\(E_\kappa\le D_L\). Appendix~\ref{sec:main-proof} gives the complete proof, including
the two cutoff regimes and the admissibility of the smoothing scale.


% Source: sharpness_and_discussion
\section{Sharpness and remaining questions}\label{opt:section}
The dependence on sample size cannot improve beyond \(n^{-1/2}\)
under the present assumptions. A normalized Rademacher sum has an atom
of order \(n^{-1/2}\), while the Gaussian probability of an interval
shrinking to this atom tends to zero. This obstruction already occurs
in one dimension and survives zero skewness.

Dimension introduces another obstruction, already visible when the
coordinates are independent.
\begin{theorem}[Lower scale]\label{opt:skew-lower}
Suppose \(p_n\to\infty\), \(B_n\ge2\), and
\(r_n=B_n(\log p_n)^{3/2}/\sqrt n\to0\). There are i.i.d.
centered observations with independent coordinates, covariance
\(I_{p_n}\), and envelope \(B_n\), for which
\begin{equation}
 \liminf_n\frac{\rh(W_n,Z_n)}{r_n}\ge\frac{\sqrt2}{4e}.
 \label{opt:lower-constant}
\end{equation}
Here \(W_n=n^{-1/2}\sum_iX_{n,i}\),
\(Z_n\sim N(0,I_{p_n})\), and the envelope is attained.
\end{theorem}
The example uses the centered, unit-variance two-point law
\[
 \Pp(X=B_n)=\frac1{B_n^2+1},\qquad
 \Pp(X=-B_n^{-1})=\frac{B_n^2}{B_n^2+1}.
\]
At a normal maximum threshold of order \(\sqrt{\log p_n}\), its
skewness changes each relative tail by order \(r_n\). Taking the
product of the coordinate probabilities preserves this discrepancy.
The Supplement proves a moderate-tail estimate uniform in the envelope
and includes the lattice corrections. This extends the regime of the
lower construction in~\cite{CCK} without changing that construction.

These lower examples have correlation gap one. To express uniformity,
let \(M_n(c_0;B,p)\) be the supremum of rectangle error over the
independent centered observations with envelope at most \(B\), unit
average marginal variances, and correlation gap at least \(c_0\).
Write \(M_n^0(c_0;B,p)\) for the same supremum restricted to
aggregate zero third-moment tensor.
For fixed \(c_0\in(0,1]\), Theorem~\ref{thm:optimized} and the lower
example give
\begin{equation}
 M_n(c_0;B_n,p_n)\asymp_{c_0,K}
 \frac{B_n(\log p_n)^{3/2}}{\sqrt n}
 \label{opt:matching-minimax}
\end{equation}
whenever \(p_n\to\infty\), \(B_n\ge2\), the displayed rate tends
to zero, and \(1+\log_+(\sqrt n/B_n)\le K\log(ep_n)\).
This includes polynomial dimension and growing envelopes subject to
the vanishing-rate condition.

\begin{theorem}[Uniform consistency at a fixed gap]
\label{opt:consistency}
For fixed \(c_0\in(0,1]\), arbitrary dimension sequences \(p_n\ge1\),
and envelopes \(B_n\ge2\),
\begin{equation}
 M_n(c_0;B_n,p_n)\longrightarrow0
 \quad\Longleftrightarrow\quad
 \frac{B_n\{\log(ep_n)\}^{3/2}}{\sqrt n}\longrightarrow0.
 \label{opt:consistency-iff}
\end{equation}
For each fixed \(B\ge1\),
\[
 M_n^0(c_0;B,p_n)\longrightarrow0
 \quad\Longleftrightarrow\quad \log p_n=o(\sqrt n).
\]
\end{theorem}

For the zero-skewness assertion, independent Rademacher coordinates
locate the critical dimension at
\(\log p\asymp\sqrt n\). If \(\log p_n\sim a\sqrt n\), their
maximum has a limiting discrepancy
\(e^{-\exp(-a^2/3)}-e^{-1}>0\) at a threshold chosen so that the
Gaussian upper-tail count tends to one. The fourth-order tail
calculation and the necessity argument for
arbitrary dimension sequences are in the Supplement. A sharp threshold
does not establish a sharp error rate: in the range
\(n^{1/3}\ll L\ll\sqrt n\), known lower bounds have scale
\(L^2/n\), and the fixed-gap zero-skewness upper bound is larger.
The lattice obstruction must also be retained in any improved
unsmoothed error bound.

For unrestricted covariance the upper bound is optimal up to
logarithmic factors. Where the lower scale \(\delta L^{3/2}\) is
attained, the remaining upper--lower gap is at most
\(\sqrt{D_L}\). Neither the necessity of this factor nor the optimal
dependence on a vanishing correlation gap is known. The companion
investigations~\cite{Cancellation} remove the factor when covariance
rank is at most two, uniformly over the number and angular separation
of the coordinate faces. They do not establish the corresponding
growing-rank result or the optimal logarithmic power within the planar
class. The transport identities and rank-induction tests are developed
in~\cite{Cancellation}.

Exactly repeated or opposite coordinates can be removed without
changing rectangle distance. The Supplement gives this reduction and
the rank-one minimax rate \(C\min\{1,B/\sqrt n\}\). These useful
special cases do not change the arbitrary-rank proof.

The tail extension in Corollary~\ref{cor:subexponential} complements the
locally nondegenerate companion paper~\cite{Local}. Its sub-exponential
bound is sharper in logarithms when all principal covariance submatrices
of order at most four are uniformly nondegenerate. Here the sample-size
power remains \(n^{-1/2}\), up to logarithms, with no such covariance
restriction. Appendix~\ref{app:subexponential} compares the assumptions and rates explicitly.


% Source: references
\begin{thebibliography}{99}
\bibitem{Cancellation}
\emph{Cancellation and rank reduction in rectangle CLTs:
transport identities, planar bounds, and rank-three tests}.
Companion research note, 20 October 2026.

\bibitem{AL}
P.~M. Aronow and P. Lopatto (2026).
\emph{Optimal central limit theorem for bounded random variables in high dimensions}.
Manuscript, arXiv:2609.07067v1.
\url{https://arxiv.org/abs/2609.07067v1}.

\bibitem{Gumbel}
\emph{The rational rectangle smoother and the Gumbel trick}.
Companion note.
\nolinkurl{Rational_Smoother_and_Gumbel_Trick.pdf}.

\bibitem{Kane}
D. Kane (2014).
The average sensitivity of an intersection of half spaces.
\emph{Research in the Mathematical Sciences} \textbf{1}, Article 13.
\href{https://doi.org/10.1186/s40687-014-0013-6}{doi:10.1186/s40687-014-0013-6}.

\bibitem{CCK}
V. Chernozhukov, D. Chetverikov, and Y. Koike (2023).
\emph{Nearly optimal central limit theorem and bootstrap approximations
in high dimensions}.
The Annals of Applied Probability \textbf{33}(3), 2374--2425.
\href{https://doi.org/10.1214/22-AAP1870}{doi:10.1214/22-AAP1870}.
Working version: arXiv:2012.09513v2, May 12, 2021.
\url{https://arxiv.org/abs/2012.09513v2}.

\bibitem{Local}
V. Chernozhukov, D. Chetverikov, K. Kato, and Y. Koike,
with AI collaborators (2026).
\emph{High-dimensional Gaussian approximation over rectangles under
local nondegeneracy}. Companion working manuscript supplied
20 October 2026.

\bibitem{CCK2013}
V. Chernozhukov, D. Chetverikov, and K. Kato (2013).
\emph{Gaussian approximations and multiplier bootstrap for maxima of
sums of high-dimensional random vectors}.
The Annals of Statistics \textbf{41}(6), 2786--2819.
\href{https://doi.org/10.1214/13-AOS1161}{doi:10.1214/13-AOS1161}.

\bibitem{CCK2017}
V. Chernozhukov, D. Chetverikov, and K. Kato (2017).
\emph{Central limit theorems and bootstrap in high dimensions}.
The Annals of Probability \textbf{45}(4), 2309--2352.
\href{https://doi.org/10.1214/16-AOP1113}{doi:10.1214/16-AOP1113}.

\bibitem{CCKK2022}
V. Chernozhukov, D. Chetverikov, K. Kato, and Y. Koike (2022).
\emph{Improved central limit theorem and bootstrap approximations in
high dimensions}. The Annals of Statistics \textbf{50}(5), 2562--2586.
\href{https://doi.org/10.1214/22-AOS2193}{doi:10.1214/22-AOS2193}.

\bibitem{GaoZhang}
Z. Gao and W. Zhang (2026).
\emph{Cubic-root Gaussian approximation under unrestricted covariance}.
arXiv:2608.30221v1, 31 August 2026.
\url{https://arxiv.org/abs/2608.30221v1}.

\bibitem{FKLZ}
X. Fang, Y. Koike, S.-H. Liu, and Y.-K. Zhao (2026).
\emph{High-dimensional central limit theorems by Stein's method in
the degenerate case}. The Annals of Applied Probability, to appear.
\href{https://research.cuhk.edu.hk/en/publications/high-dimensional-central-limit-theorems-by-steins-method-in-the-d/}{Accepted manuscript (2026)}.

\end{thebibliography}



\clearpage
\appendix
\numberwithin{equation}{section}
\numberwithin{theorem}{section}
\numberwithin{lemma}{section}
\numberwithin{proposition}{section}
\numberwithin{corollary}{section}
\numberwithin{remark}{section}
% Source: smoothing_proofs
\section{Smoothing a rectangle}\label{appx-sec:smoothing}

We replace the indicator of a rectangle by a smooth function whose error
and derivatives are concentrated near the boundary.  The same
concentration estimate will control both quantities.  We state that
estimate first, then construct the smoothing function and bound the
error it introduces at \(W\) and \(Z\).

\subsection{Concentration near the boundary}

For a rectangle \(A=\prod_j[\ell_j,u_j]\), let
\(A^a=\prod_j[\ell_j-a,u_j+a]\); a reversed interval is empty.
The same notation for a fixed system of halfspaces increases all its
thresholds by \(a\).

\begin{lemma}[Rectangle concentration]\label{appx-lem:concentration}
Let \(S=\sum_i\zeta_i\), where the summands are independent, centered,
and bounded by \(\eta\) in maximum norm. Suppose
\(\Var(S_j)=v>0\) for every coordinate. For every independent random
shift \(R\), every rectangle \(A\), and \(a>0\),
\begin{equation}
 \Pp(S+R\in A^a\setminus A^{-a})
 \le C\min\{1,(a+\eta)\sqrt{\log(ep)/v}\}.
 \label{appx-eq:concentration}
\end{equation}
For a centered Gaussian base with coordinate variances \(v\), the same
bound holds with \(\eta=0\), including singular covariance matrices.
\end{lemma}

A proof is given in Appendix~\ref{appx-app:concentration}.
We also apply the lemma to the signed face vector
\((\nu_\alpha\cdot x)_{\alpha\le m}\). Its coordinates retain
the same variance and summand bound. The intersection
\(\{\nu_\alpha\cdot x\le a_\alpha\}_{\alpha\le m}\) is the inverse
image of the lower orthant \(\prod_{\alpha\le m}(-\infty,a_\alpha]\)
in \(\R^m\); comparing thresholds increased and decreased by \(a\)
gives the shell in the concentration estimate. The orthant remains nonempty even when
its inverse image is empty, so inconsistent threshold systems are
covered. Since \(m\le2p\), \(\log(em)\le CL\).

\subsection{A selector and its derivatives}

We encode the face violations by positive exponential weights.  Their
sum will give a smooth approximation to the rectangle indicator.  Write
a finite rectangle as
\[
 A=\{x:\nu_\alpha\cdot x\le a_\alpha,\ \alpha\in\cF\},
\]
where the normals are distinct members of
\(\{\pm e_1,\ldots,\pm e_p\}\). Put \(m=|\cF|\le2p\), and define
\begin{equation}
 \begin{aligned}
 w_\alpha(x)&=e^{(\nu_\alpha\cdot x-a_\alpha)/h},
 &T(x)&=\sum_\alpha w_\alpha(x),\\
 F_{A,h}(x)&=\frac1{1+T(x)},
 &\Psi_{A,h}(x)&=h\log(1+T(x)).
 \end{aligned}
 \label{appx-eq:selector}
\end{equation}
If \(m=0\), the selector is identically one and all comparisons are
trivial; the analytic arguments below take \(m\ge1\).
We usually suppress \(A,h\) and write
\[
 \pi_0=F,\qquad \pi_\alpha=\frac{w_\alpha}{1+T},
 \qquad \theta=1-F,\qquad v=\sum_\alpha\pi_\alpha\nu_\alpha,
 \qquad \chi=h\,\tr D^2\Psi.
\]
The rational form of \(F\) has a useful probabilistic interpretation:
it is an average of rectangle indicators.  If
\(E_0,(E_\alpha)_\alpha\) are independent unit exponentials, then
\begin{equation}
 F_{A,h}(x)=\E\ind_{A_E}(x),\qquad
 A_E=\{x:\nu_\alpha\cdot x\le
 a_\alpha+h\log(E_\alpha/E_0)\text{ for every }\alpha\}.
 \label{appx-eq:exp-mixture}
\end{equation}
Indeed, conditioning on \(E_0\), the probability on the right is
\(e^{-E_0T(x)}\), whose expectation is \((1+T(x))^{-1}\).
Equivalently, \(G_\beta=-\log E_\beta\) are independent standard
Gumbel variables, so
\[
 F_{A,h}(x)=\Pp\!\left\{hG_0\ge
 \max_\alpha(\nu_\alpha\cdot x-a_\alpha+hG_\alpha)\right\}.
\]
The auxiliary option has deterministic score zero. The face displacements
\(h(G_0-G_\alpha)\) have logistic marginals and are dependent through
\(G_0\); independent logistic displacements give a different selector.
This is the Gumbel interpretation of \cite{Gumbel}.

We will use derivatives of orders three and four.  It is convenient to
record their common formula: for vectors \(z_1,\ldots,z_q\), the chain
rule gives
\begin{equation}
 \begin{split}
 D^qF(x)[z_1,\ldots,z_q]
 ={}&h^{-q}\sum_{\pi\in\cP_q}(-1)^{|\pi|}|\pi|!
 \sum_{\alpha_1,\ldots,\alpha_{|\pi|}}
 F(x)\prod_{b=1}^{|\pi|}\pi_{\alpha_b}(x)\\
 &\hspace{3em}\times
 \prod_{b=1}^{|\pi|}\prod_{a\in B_b}
 (\nu_{\alpha_b}\cdot z_a),
 \end{split}
 \label{appx-eq:partition-derivative}
\end{equation}
where \(\cP_q\) is the set of partitions of \(\{1,\ldots,q\}\),
and \(B_1,\ldots,B_{|\pi|}\) are the blocks in any fixed ordering.
Ordered partitions can be encoded by a permutation and cuts between
successive elements, so
\(\sum_{\pi\in\cP_q}|\pi|!\le q!2^{q-1}\).
In particular,
\begin{equation}
 |D^qF(x)[z_1,\ldots,z_q]|
 \le C^q q!h^{-q}\prod_{a=1}^q
 \max_\alpha|\nu_\alpha\cdot z_a|.
 \label{appx-eq:all-orders}
\end{equation}

The uniform bound \eqref{appx-eq:all-orders} does not yet record that the
derivatives are small away from the boundary.  For that purpose, let
\(I=(i_1,\ldots,i_q)\) be an ordered coordinate tuple and write
\(\partial_I=\partial_{i_1}\cdots\partial_{i_q}\).
The same positive majorant, using \(\norm{\nu_\alpha}_1=1\), gives
\begin{equation}
 \sum_I\sup_{\norm z_\infty\le r}|\partial_IF(x+z)|
 \le C_qh^{-q}F(x)(1-F(x)),\qquad 0\le r\le h.
 \label{appx-eq:local-envelope}
\end{equation}
To see the local stability, each \(w_\alpha\) changes by a factor in
\([e^{-r/h},e^{r/h}]\), and \(1+T\) obeys the same comparison.
After summing over coordinate tuples and face indices, the terms with
\(k\) blocks are bounded by a constant times
\(h^{-q}T^k/(1+T)^{k+1}\le h^{-q}F(1-F)\).

\begin{lemma}[Averaging the derivative envelope]\label{appx-lem:envelope}
Suppose a random vector \(U\) has an independent component that is
either a sum of independent centered vectors, with summand bound
\(2\delta\) and equal coordinate variances at least \(1/2\), or a
centered Gaussian vector with equal coordinate variances at least
\(1/2\).
Uniformly over thresholds, for \(h\ge4\delta\),
\begin{equation}
 \E F(U)(1-F(U))\le Ch\sqrt L.
 \label{appx-eq:small-envelope}
\end{equation}
Consequently, for fixed \(q\) and \(h\ge\max\{4\delta,r\}\),
\begin{equation}
 \E\sum_I\sup_{\norm z_\infty\le r}|\partial_IF(U+z)|
 \le C_q h^{1-q}\sqrt L.
 \label{appx-eq:averaged-derivatives}
\end{equation}
\end{lemma}

\begin{proof}
Shift every face threshold by the same amount \(s\), and write
\(F_s=(1+e^{-s/h}T)^{-1}\) for the resulting selector. Since
\(\partial_sF_s=F_s(1-F_s)/h\), and
\(F_s(1-F_s)\ge e^{-1}F(1-F)\) for \(|s|\le h\),
\[
 F(1-F)\le\frac e2(F_h-F_{-h}).
\]
In the mixture representation~\eqref{appx-eq:exp-mixture},
\(F_h-F_{-h}\) is the average of
\(\ind_{A_E^h}-\ind_{A_E^{-h}}\), with shifts interpreted in the
signed face coordinates. Conditional on the exponentials and the
independent shift, Lemma~\ref{appx-lem:concentration} bounds the expectation
by \(C(h+\delta)\sqrt L\): the base has equal coordinate variances
at least \(1/2\) and summand bound at most \(2\delta\), or is
Gaussian. Since \(h\ge4\delta\), this proves
\eqref{appx-eq:small-envelope}. Combining it with
\eqref{appx-eq:local-envelope} proves \eqref{appx-eq:averaged-derivatives}.
\end{proof}

\subsection{The error at the endpoints}

The mixture representation gives a bound for every covariance matrix.
A second argument uses the correlation gap and can be sharper.  We will
choose the better of the two bounds after the smoothing scale has been
fixed.

\begin{lemma}[Unrestricted endpoint]\label{appx-lem:endpoint-free}
Under \eqref{eq:setup}, for \(h\ge\delta\),
\begin{equation}
 \E|F_{A,h}(W)-\ind_A(W)|
 +\E|F_{A,h}(Z)-\ind_A(Z)|
 \le ChL^{3/2}.
 \label{appx-eq:endpoint-free}
\end{equation}
\end{lemma}

\begin{proof}
In \eqref{appx-eq:exp-mixture}, let
\(R_h=h\max_\alpha|\log(E_\alpha/E_0)|\).
Then \(A^{-R_h}\subseteq A_E\subseteq A^{R_h}\), and
\(\E R_h\le Ch\log(em)\le ChL\).
The last assertion follows by integrating the exponential tails of
\(-\log E_\alpha\), and the faster upper tails of \(\log E_\alpha\).
For fixed exponentials, the indicators can disagree only in
\(A^{R_h}\setminus A^{-R_h}\). Jensen's inequality moves the absolute
value inside the exponential average. Conditional on the exponentials,
Lemma~\ref{appx-lem:concentration} bounds the errors at \(W\) and \(Z\)
by \(C(R_h+\delta)\sqrt L\) and \(CR_h\sqrt L\), respectively.
Averaging the radius gives \(C(hL+\delta)\sqrt L\) and
\(ChL\sqrt L\).
Use \(h\ge\delta\) and \(L\ge1\).
\end{proof}

\begin{lemma}[Curvature endpoint]\label{appx-lem:endpoint-gap}
Under \eqref{eq:setup}, if \(c>0\) and \(h\ge\delta\), then
\begin{equation}
 \E|F_{A,h}(W)-\ind_A(W)|
 +\E|F_{A,h}(Z)-\ind_A(Z)|
 \le \frac{Ch\sqrt L}{c}.
 \label{appx-eq:endpoint-gap}
\end{equation}
\end{lemma}

\begin{proof}
The curvature of \(\Psi\) controls the error of \(F\) pointwise.
Direct differentiation gives
\begin{equation}
 D^2\Psi=\frac1h\left(
 \sum_\alpha\pi_\alpha\nu_\alpha\nu_\alpha^\top-vv^\top\right),
 \qquad \norm v_1\le\theta,\qquad
 \chi=\theta-\norm v_2^2.
 \label{appx-eq:hessian}
\end{equation}
Distinct signed coordinate normals have nonpositive inner products.
Thus, with \(\pi_* =\max_\alpha\pi_\alpha\),
\(\chi\ge\theta(1-\theta)\) and
\(\chi\ge\theta(1-\pi_*)\).
On \(A\), every \(w_\alpha\le1\), so \(\pi_*\le1/2\).
Off \(A\), \(T>1\), so \(\theta>1/2\). On \(A\), the second
lower bound gives \(\chi\ge\theta/2\); off \(A\), the first gives
\(\chi\ge F/2\). The indicator error is \(\theta\) and \(F\)
in the two cases, respectively, so
\begin{equation}
 |F-\ind_A|\le2\chi.
 \label{appx-eq:curvature-error}
\end{equation}
To bring the covariance into this estimate, observe that
\[
 \sum_{j\ne k}|\partial_{jk}\Psi|
 =h^{-1}(\norm v_1^2-\norm v_2^2)
 \le\tr D^2\Psi,
\]
and since \(\Sigma_{jj}=1\) and \(|\Sigma_{jk}|\le1-c\) for
\(j\ne k\), the off-diagonal contribution can remove at most
\((1-c)\tr D^2\Psi\). Therefore
\begin{equation}
 \tr(\Sigma D^2\Psi)\ge c\,\tr D^2\Psi.
 \label{appx-eq:gap-contraction}
\end{equation}

We next compare the Hessian at a point and at a nearby point.  Include
the null face \(\nu_0=0\).  The positive representation
\[
 D^2\Psi=\frac1{2h}\sum_{\alpha,\beta\in\{0\}\cup\cF}
 \pi_\alpha\pi_\beta
 (\nu_\alpha-\nu_\beta)(\nu_\alpha-\nu_\beta)^\top
\]
implies
\(e^{-4r/h}D^2\Psi(x)\preceq D^2\Psi(x+z)
\preceq e^{4r/h}D^2\Psi(x)\) when \(\norm z_\infty\le r\).
Set \(W_i=W-\xi_i\). Centering and independence let us subtract
\(\nabla\Psi(W_i)\) before integrating the Hessian along the segment
from \(W_i\) to \(W\). First compare this Hessian with
\(D^2\Psi(W_i)\). Independence then replaces
\(\xi_i\xi_i^\top\) by \(\E\xi_i\xi_i^\top\). A second comparison
replaces \(D^2\Psi(W_i)\) by \(D^2\Psi(W)\), and summing the row
covariances gives \(\Sigma\). Each lower comparison contributes a factor
\(e^{-4\delta/h}\), so
\[
 \begin{split}
 \E[W\cdot\nabla\Psi(W)]
 &=\sum_i\E\int_0^1
 \xi_i^\top D^2\Psi(W_i+s\xi_i)\xi_i\dd s\\
 &\ge e^{-8\delta/h}\E\tr(\Sigma D^2\Psi(W)).
 \end{split}
\]
Since \(\norm{\nabla\Psi}_1\le1\), the left side is at most
\(\E\norm W_\infty\), which Bernstein's inequality bounds by
\(C(\sqrt L+\delta L)\). When \(\delta\sqrt L\le1\), this is
\(C\sqrt L\), and \(h\ge\delta\) bounds \(e^{8\delta/h}\)
absolutely. Combining with \eqref{appx-eq:gap-contraction} gives
\(\E\chi(W)\le Ch\sqrt L/c\).
If \(\delta\sqrt L>1\), the same conclusion follows from
\(\chi\le1\) and \(h\ge\delta\).
For \(Z\), Gaussian integration by parts gives
\[
 \E[Z\cdot\nabla\Psi(Z)]
 =\E\tr(\Sigma D^2\Psi(Z)),\qquad
 \E\norm Z_\infty\le C\sqrt L.
\]
This identity remains valid for singular \(\Sigma\), by restricting to
its support. Use \eqref{appx-eq:curvature-error} to finish.
\end{proof}


Combining the two endpoint bounds gives, with the
parameter \(\kappa\) defined in Section~\ref{sec:results},
\begin{equation}
 \E|F_{A,h}(W)-\ind_A(W)|+
 \E|F_{A,h}(Z)-\ind_A(Z)|
 \le Ch\sqrt L\,\kappa,\qquad h\ge\delta.
 \label{appx-eq:endpoint-min}
\end{equation}


% Source: covariance_support
\section{Gaussian estimates and interpolation}\label{appx-sec:support}

We now compare the expectations of the selector at \(W\) and \(Z\).
The comparison follows a path that preserves covariance and gradually
replaces the bounded summands by a Gaussian vector.  Its derivative
contains third and fourth derivatives of the selector.  The difficulty is to sum over rows without paying for the support
dimension. We first prove a surface estimate, turn it into a matrix
bound, and use the row covariance budget to control the derivatives.
Subsection~\ref{appx-an:sec:path} then introduces the path and applies these
estimates.

\subsection{Coordinates on the covariance support}

Working on the support of \(\Sigma\) lets us use integration by parts
against a nondegenerate standard Gaussian, even when \(\Sigma\) is
singular.
Factor \(\Sigma=BB^\top\), where \(B\) is \(p\times r\) of full
column rank and \(r=\operatorname{rank}(\Sigma)\).
Each \(\xi_i\) belongs almost surely to the range of \(B\): every
null direction of \(\Sigma\) has zero variance for each centered row.
We may therefore write
\begin{equation}
 \xi_i=B\zeta_i^0,\qquad
 \sum_i\E \zeta_i^0(\zeta_i^0)^\top=I_r,\qquad Z\stackrel d=BG,
 \qquad G\sim N(0,I_r).
 \label{appx-eq:support}
\end{equation}
Set \(u_\alpha=B^\top\nu_\alpha\). Unit marginal variances imply
\(\norm{u_\alpha}=1\), and
\begin{equation}
 \norm{\zeta_i^0}_{\mathcal U}
 :=\max_\alpha|u_\alpha\cdot \zeta_i^0|\le\delta.
 \label{appx-eq:projected-bound}
\end{equation}
Thus every face projection of a support row is the corresponding signed
coordinate of the original row. Passing to the support changes the Gaussian
dimension, but preserves the directional bounds used by the selector.
The inverse image of a rectangle is a polyhedron with these normals.
For the rest of this section, \(F\) denotes the restriction
\(x\mapsto F_{A,h}(Bx)\), and its weights are restricted in the same way.
All directional derivative identities above hold with \(u_\alpha\)
in place of \(\nu_\alpha\).



% Source: Gaussian_and_interpolation_proofs
\subsection{An averaged Gaussian surface estimate}
\label{appx-an:sec:surface}

The first estimate concerns Gaussian surface integrals of randomly
translated polyhedra.  Taking the absolute value after the surface
integral preserves a cancellation that will be needed in the matrix
estimate of Subsection~\ref{appx-an:sec:matrix}. Nearly parallel facets
require care: their separate ridge terms contain inverse angles. Pairing
those terms before estimating them will remove this loss.

All Gaussian integrations take place in the covariance
support \(\R^r\) introduced in \eqref{appx-eq:support}.  Thus the face normals
\(u_1,\ldots,u_m\) have Euclidean norm one, \(m\le2p\), and
\(\norm{x}_U=\max_\alpha|u_\alpha\cdot x|\).
For \(a\in\R^m\), write
\[
 Q(a)=\{x:u_\alpha\cdot x\le a_\alpha,\ \alpha\le m\},
 \qquad Q(a)^t=Q(a+t\mathbf1).
\]
We allow empty systems and inconsistent systems.  The random
translation \(V_0\) will have the concentration property
\begin{equation}
 \Pp\{V_0\in Q(a)^t\setminus Q(a)^{-t}\}
 \le C_0(t+\delta)\sqrt L,
 \qquad a\in\R^m,\quad t>0.
 \label{appx-an:concentration}
\end{equation}
The same property then holds for subsystems, by sending unused
thresholds to infinity, and for deterministic translates of \(V_0\).
Lemma~\ref{appx-lem:concentration} supplies this property below with an
absolute constant \(C_0\).

Let \(\Ga\) and \(\ph\) be standard Gaussian measure and density
in \(\R^r\), and put
\(\Hh(g)=\norm{g}^2-r\).  For a proper full-dimensional polyhedron
\(Q\), define
\[
 \surf(Q)=\int_{\partial Q}\Hh(g)\ph(g)\dd\mathcal H^{r-1}(g),
\]
where actual facets are counted once.  Set \(\surf(Q)=0\) for the
empty set and the full space.  Lower-dimensional transition levels
will occur only at exceptional values of a common threshold shift and
will be handled by first variation.

\begin{lemma}[Averaged scalar surface]\label{appx-an:lem:surface}
Suppose \eqref{appx-an:concentration} holds, let \(G\sim N(0,I_r)\) be
independent of \(V_0\), and let \(0<\sigma\le1\).  Uniformly over
empty or full-dimensional target polyhedra \(Q\) with the specified normals,
\begin{equation}
 \E\left|\surf\bigl((Q-V_0)/\sigma\bigr)\right|
 \le C\bigl(\sigma L^{3/2}+\delta L^2\bigr).
 \label{appx-an:surface-bound}
\end{equation}
Here and below, a target which is lower dimensional may be omitted
when integrating over an absolutely continuous common threshold shift.
In particular, if \(\sigma\ge\delta\sqrt L\), the right-hand side of
\eqref{appx-an:surface-bound} is at most \(C\sigma L^{3/2}\).
\end{lemma}

\begin{proof}
Condition first on \(V_0\), and put \(Q'=(Q-V_0)/\sigma\).
Assume that \(Q'\) is proper and full dimensional; remove redundant
inequalities and repeated facet descriptions.  Denote a facet's
offset by \(a'_j\) and its Gaussian surface mass by \(s_j\).
For a ridge shared by facets \(j,k\), let \(z_{jk}\) be the constant
orthogonal projection of its points onto
\(\operatorname{span}(u_j,u_k)\), let \(\rho_{jk}\) be its Gaussian
codimension-two surface mass, and put
\(\alpha_{jk}=\arccos(u_j\cdot u_k)\in(0,\pi)\).
Tangent Gaussian divergence on every facet gives the identity
\begin{equation}
 \surf(Q')=
 \sum_j\bigl((a'_j)^2-1\bigr)s_j
 -\sum_{\{j,k\}}(d_{k\mid j}+d_{j\mid k})\rho_{jk},
 \label{appx-an:surface-identity}
\end{equation}
where
\[
 d_{k\mid j}=
 \frac{a'_k-(u_j\cdot u_k)a'_j}{\sin\alpha_{jk}},
 \qquad
 d_{k\mid j}+d_{j\mid k}
 =2\sin(\alpha_{jk}/2)
    \frac{u_j+u_k}{\norm{u_j+u_k}}\cdot z_{jk}.
\]
Indeed, on facet \(j\), write \(g=a'_ju_j+y\), with \(y\perp u_j\),
and use
\[
 (\norm{y}^2-(r-1))\varphi_{r-1}(y)
 =-\operatorname{div}_{u_j^\perp}
                  \{y\varphi_{r-1}(y)\}.
\]
The boundary term on the ridge with facet \(k\) is
\(-d_{k\mid j}\rho_{jk}\).  Pairing its two contributions gives
\eqref{appx-an:surface-identity}.  Truncation justifies this calculation
for unbounded facets, since polynomial Gaussian terms at infinity
vanish.  When \(r=1\), there are no ridge terms and the same identity
is interpreted as a sum over endpoints.

We estimate the facet and ridge terms by the Gaussian mass of short
outward tubes.  Pairing the two ridge contributions before estimating
them makes the bound uniform in the angle between the facets.  Fix
\(R\ge2\) and
set \(\varepsilon_R=(1+R)^{-1}\).  Above a facet with
\(|a'_j|\le R\), its outward normal tube of lengths
\(0<t<\varepsilon_R\) has Gaussian mass \(q_j\) satisfying
\[
 q_j\ge c\varepsilon_Rs_j,
 \qquad (1+(a'_j)^2)s_j\le CR^3q_j.
\]
The density ratio at a facet point \(x\) displaced by \(tu_j\) is
\(\exp(-ta'_j-t^2/2)\), which is bounded below by an absolute
constant on this tube.  For a ridge with \(\norm{z_{jk}}\le R\),
use displacements of norm less than \(\varepsilon_R\) in its
two-dimensional outward normal cone.  That cone is a sector of angle
\(\alpha_{jk}\).  For a displacement \(y\) in it,
\(x\cdot y=z_{jk}\cdot y\), so the same density-ratio bound gives
\[
 q_{jk}\ge c\alpha_{jk}\varepsilon_R^2\rho_{jk}.
\]
Since
\( |d_{k\mid j}+d_{j\mid k}|\le C\alpha_{jk}R\), it follows that
\[
 |d_{k\mid j}+d_{j\mid k}|\rho_{jk}\le CR^3q_{jk}.
\]
All these tubes are mutually disjoint up to null sets: their points
have distinct metric projections onto \(Q'\).  They lie in
\((Q')^{\varepsilon_R}\setminus Q'\).  These assertions use the
actual faces of the polyhedron and remain valid if the original
threshold system contained redundant inequalities.

To sum these local estimates, divide the facet and ridge offsets into
dyadic shells.  Let \(R_\ell=2^{\ell+1}\),
\(\ell\ge0\); the zeroth shell contains facet offsets of absolute
value at most \(2\) and ridge offsets of norm at most \(2\), and shell
\(\ell\ge1\) contains the corresponding offsets in
\((R_\ell/2,R_\ell]\).  Let \(T_\ell(V_0)\) be the union of the
associated tubes, all of radius \((1+R_\ell)^{-1}\).  Then
\begin{equation}
 |\surf(Q')|\le C\sum_{\ell\ge0}R_\ell^3
                         \Ga(T_\ell(V_0)).
 \label{appx-an:dyadic-tubes}
\end{equation}
For linearly independent pairs, let \(\Pi_{jk}\) be orthogonal
projection onto \(\operatorname{span}(u_j,u_k)\), and set
\[
 K(G)=\max\left\{\max_j|u_j\cdot G|,
   \max_{j,k:\,u_j,u_k\text{ independent}}\norm{\Pi_{jk}G}\right\}.
\]
Omit the second maximum when it is empty.  If
\(G\in T_\ell(V_0)\) and \(\ell\ge1\), then
\(K(G)\ge R_\ell/2-1\): its appropriate facet coordinate or ridge
projection differs from the shell offset by at most one.
Also, for every \(\ell\),
\[
 G\in T_\ell(V_0)
 \quad\Longrightarrow\quad
 V_0+\sigma G\in Q^{\sigma/(1+R_\ell)}\setminus Q.
\]
Condition on \(G\), apply \eqref{appx-an:concentration} to the last event,
and retain the necessary event involving \(K(G)\).  This gives
\[
 \E\Ga(T_\ell(V_0))
 \le C(\sigma/R_\ell+\delta)\sqrt L\,\beta_\ell,
\]
where \(\beta_0=1\) and
\(\beta_\ell=\Pp\{K(G)\ge R_\ell/2-1\}\) for \(\ell\ge1\).
There are at most \(m+m^2\) one- or two-dimensional projections in
this maximum, whence Gaussian tails imply
\[
 \beta_\ell\le\min\{1,Cm^2e^{-cR_\ell^2}\},\qquad \ell\ge1.
\]
Splitting the dyadic sum at \(R_\ell\) of order
\(\sqrt{\log(em)}\) yields, for \(q=2,3\),
\[
 \sum_{\ell\ge0}R_\ell^q\beta_\ell
 \le C_q\{\log(em)\}^{q/2}\le C_qL^{q/2}.
\]
Taking expectations in \eqref{appx-an:dyadic-tubes} and using Tonelli's
theorem proves \eqref{appx-an:surface-bound}.
\end{proof}

\paragraph{First variation and common shifts.}
In the matrix estimate we will differentiate a random common shift of
all thresholds. The following formula justifies this operation even when
a polyhedron collapses at an exceptional shift. For fixed offsets and \(V_0\), the function
\[
 s\longmapsto
 \E_G\left[\Hh(G)\ind\left\{
   \max_j\{u_j\cdot(V_0+\sigma G)-a_j\}\le s\right\}\right]
\]
is absolutely continuous.  The inverse image of a Lebesgue-null set
under the maximum is contained in the union of its inverse images
under the individual nondegenerate Gaussian face coordinates.  Coarea
therefore gives its derivative, for almost every \(s\), as
\(\sigma^{-1}\surf((Q(a+s\mathbf1)-V_0)/\sigma)\).
To justify interchanges before applying the dimension-free estimate,
a finite bound \(Cm(r+1)/\sigma\) on its Lipschitz constant suffices.
Indeed, the integral of \(|\Hh|\) over a hyperplane at standard
Gaussian coordinate \(t\) is at most
\((t^2+2r)\varphi_1(t)\le C(r+1)\), and there are at most \(m\)
facets.  Independent absolutely continuous common threshold shifts
avoid exceptional levels and allow these first variations to be
averaged.  The temporary dimension dependence in this justification
does not enter any asserted bound.

\subsection{A matrix estimate for the selector weights}
\label{appx-an:sec:matrix}

We next turn the scalar surface estimate into a bound for a Gaussian
matrix weighted by selector derivatives.  The matrix norm is useful
because it can be paired with a covariance budget without summing over
the support dimension.

Recall the selector notation
\begin{gather*}
 t_j(x)=\frac{u_j\cdot x-a_j}{h},\qquad
 N(x)=1+\sum_je^{t_j(x)},\\
 F(x)=N(x)^{-1},\qquad
 \pi_j(x)=\frac{e^{t_j(x)}}{N(x)},\qquad \theta(x)=1-F(x).
\end{gather*}
In particular, \(\sum_j\pi_j=\theta\).  If \(h\ge C\delta\),
the exponential rectangle mixture for \(F\), coupled at common
threshold shifts \(\pm h\), and \eqref{appx-an:concentration} give
\begin{equation}
 \sup_z\E F(V_0+z)\theta(V_0+z)\le Ch\sqrt L.
 \label{appx-an:weight-concentration}
\end{equation}
The analogous one-face bound holds for every subsystem consisting
of a single inequality, with the same upper bound.
For matrices, \(\norm{\cdot}_{\mathrm{op}}\) and
\(\norm{\cdot}_*\) denote operator and nuclear norm, respectively.

A selector weight can move a Gaussian mean both normally and tangentially
to a face. The next estimate keeps the normal component as a scalar second
moment and controls the tangential component by the loss of mass relative
to a one-face weight.

\begin{lemma}[Gaussian mean relative to a single face]
\label{appx-an:lem:mean}
Let \(u\) be a unit vector, and let
\(b(g)=\beta(u\cdot g)\in[0,1]\) have positive Gaussian integral
\(p_b\).  If \(0\le w\le b\), \(p_w=\E w(G)>0\), and
\(\mu_w=\E[Gw(G)]/p_w\), then
\begin{equation}
 p_w\norm{\mu_w}^2
 \le \E[w(G)(u\cdot G)^2]
        +2p_w\log(p_b/p_w).
 \label{appx-an:mean-entropy}
\end{equation}
\end{lemma}

\begin{proof}
For \(t\perp u\), Jensen's inequality and independence of the
normal and tangent Gaussian components give
\[
 p_we^{t\cdot\mu_w}
 \le\E[w(G)e^{t\cdot G}]
 \le\E[b(G)e^{t\cdot G}]
 =p_be^{\norm{t}^2/2}.
\]
Take \(t\) equal to the projection of \(\mu_w\) onto \(u^\perp\).
The normal component satisfies
\(p_w|u\cdot\mu_w|^2\le\E[w(G)(u\cdot G)^2]\) by
Cauchy--Schwarz.  Adding the two estimates proves the result.
\end{proof}

\begin{lemma}[Matrix estimate for selector weights]
\label{appx-an:lem:matrix}
Suppose \eqref{appx-an:concentration} holds, \(h\ge C\delta\), and
\(\delta\sqrt L\le\sigma\le1\).  Put \(U=V_0+\sigma G\),
where \(G\) is independent of \(V_0\).  For every integer \(k\ge1\),
\begin{equation}
 \E_{V_0}\sum_{\alpha_1,\ldots,\alpha_k}
 \left\|\E_G\left[(GG^\top-I_r)F(U)
                \prod_{a=1}^k\pi_{\alpha_a}(U)\right]\right\|_*
 \le C^{k+1}hL^{3/2}.
 \label{appx-an:matrix-bound}
\end{equation}
\end{lemma}

\begin{proof}
Each selector weight defines a Gaussian tilt whose covariance is at
most \(I_r\). This turns the nuclear norm into a mean term and a signed
trace. We bound the means by the one-face estimate and the summed trace
by the surface estimate. Condition on \(V_0\), and write
\[
 w_\alpha(g)=F(V_0+\sigma g)
                  \prod_{a=1}^k\pi_{\alpha_a}(V_0+\sigma g),
 \quad p_\alpha=\E_Gw_\alpha(G),
 \quad q(g)=\sum_\alpha w_\alpha(g)=F(U)\theta(U)^k.
\]
Every \(w_\alpha\) is log-concave: its logarithm is a linear
function minus \((k+1)\log N(V_0+\sigma g)\).  Under its normalized
Gaussian tilt, let \(\mu_\alpha\) and \(K_\alpha\) be the mean and
covariance.  Pr\'ekopa's theorem says that
\[
 z\longmapsto\int_{\R^r}e^{-\norm{g-z}^2/2}w_\alpha(g)\dd g
\]
is log-concave.  Differentiating its logarithm at zero gives Hessian
\(K_\alpha-I_r\), so \(K_\alpha\preceq I_r\).  Thus
\begin{equation}
 \left\|\E_G[(GG^\top-I_r)w_\alpha(G)]\right\|_*
 \le2p_\alpha\norm{\mu_\alpha}^2
             -\E_G[\Hh(G)w_\alpha(G)].
 \label{appx-an:matrix-trace-reduction}
\end{equation}
To see this, write the matrix as
\(p_\alpha(K_\alpha-I_r+\mu_\alpha\mu_\alpha^\top)\), use
\(I_r-K_\alpha\succeq0\), and express its trace as
\(p_\alpha(\tr K_\alpha-r+\norm{\mu_\alpha}^2)\).

We first control the means, using the first face in each tuple as a
one-dimensional reference. For a tuple whose first entry is \(j\), put
\[
 b_j(g)=\frac{e^{t_j(V_0+\sigma g)}}
                   {(1+e^{t_j(V_0+\sigma g)})^2},
 \qquad p_{b_j}=\E_Gb_j(G).
\]
It depends only on \(u_j\cdot g\), and
\(w_\alpha\le F\pi_j\le b_j\).  Lemma~\ref{appx-an:lem:mean} therefore
implies, with \(K_1(G)=\max_j|u_j\cdot G|\),
\begin{equation}
 \sum_\alpha p_\alpha\norm{\mu_\alpha}^2
 \le\E_G[q(G)K_1(G)^2]
       +2\sum_\alpha p_\alpha\log(p_{b_{\alpha_1}}/p_\alpha).
 \label{appx-an:entropy-sum}
\end{equation}
Conditioning on \(G\), \eqref{appx-an:weight-concentration} and
\(q\le F\theta\) bound the expectation of the first term by
\(Ch\sqrt L\,\E K_1(G)^2\le ChL^{3/2}\).

For the second term, there are \(m^{k-1}\) tuples with a given first
face. The log-sum inequality lets us sum their reference masses before
estimating the entropy term. Put
\[
 \bar p=\E_{V_0,G}q(G),\qquad
 \bar B=m^{k-1}\sum_j\E_{V_0,G}b_j(G).
\]
The log-sum inequality followed by concavity of
\((x,y)\mapsto x\log(y/x)\) gives
\[
 \E_{V_0}\sum_\alpha p_\alpha
          \log(p_{b_{\alpha_1}}/p_\alpha)
 \le\bar p\log(\bar B/\bar p).
\]
We have \(\bar p\le Ch\sqrt L\).  Applying the one-face version of
\eqref{appx-an:weight-concentration} after conditioning on \(G\) gives
\(\E b_j(G)\le Ch\sqrt L\); hence
\(\bar B\le Cm^kh\sqrt L\).  Write \(\bar p=h\sqrt L\,x\),
with \(0<x\le C\).  It follows that
\[
 \bar p\log(\bar B/\bar p)
 \le h\sqrt L\,x\log(Cm^k/x)
 \le C(k+1)hL^{3/2}.
\]
The case \(\bar p=0\) follows by continuity.

We now estimate the signed trace in the sum of
\eqref{appx-an:matrix-trace-reduction}. Summing the weights first gives
\(F\theta^k\); the Gaussian score is still signed at this stage.
Shift every threshold by \(s\),
and denote the resulting functions by \(F_s,\theta_s\).  Then
\[
 F\theta^k=-\frac hk\left.\partial_s\theta_s^k\right|_{s=0}.
\]
For an integer \(a\ge1\), let \(T_a\sim\operatorname{Gamma}(a,1)\)
and let \(E_j\) be independent unit exponentials. These mixture
variables are independent of one another and of \(V_0,G\). The Gamma Laplace
transform gives the exact representation
\begin{equation}
 F_s(x)^a=
 \Pp\{u_j\cdot x\le a_j+s+h\log(E_j/T_a),\ j\le m\}.
 \label{appx-an:gamma-mixture}
\end{equation}
Expand \(\theta_s^k=\sum_{a=0}^k(-1)^a\binom ka F_s^a\).
The sum of the absolute coefficients is \(2^k\), and the term
\(a=0\) has zero derivative.  First variation,
\eqref{appx-an:gamma-mixture}, and Lemma~\ref{appx-an:lem:surface} yield
\[
 \left|\E_{V_0,G}\Hh(G)F(U)\theta(U)^k\right|
 \le\frac{2^kh}{k\sigma}
   \sup_Q\E_{V_0}|\surf((Q-V_0)/\sigma)|
 \le C2^khL^{3/2}.
\]
The common shift \(-h\log T_a\) has an absolutely continuous law;
the first-variation discussion above therefore justifies the mixture
and its differentiation, including collapsed transition levels.
Combining this estimate with \eqref{appx-an:entropy-sum} and
\eqref{appx-an:matrix-trace-reduction} proves \eqref{appx-an:matrix-bound}.
\end{proof}

\subsection{Putting translated derivatives over a common base}
\label{appx-an:sec:restoration}

In the interpolation, deleting a row produces a different base point
for each summand.  To apply the covariance budget, we put these
summands over a common base.  The rational selector lets
us do this by an exact convergent expansion of its weights.

\begin{lemma}[Paired derivatives at translated bases]
\label{appx-an:lem:restoration}
Let a sigma-field \(\mathcal H\) determine \(V_0\), let
\(G\sim N(0,I_r)\) be independent of \(\mathcal H\), and put
\(U=V_0+\sigma G\).  Assume the hypotheses of
Lemma~\ref{appx-an:lem:matrix}.  For \(s\in\{1,2\}\), let
\((Z_i,Z_{i,1},\ldots,Z_{i,s},v_i)\) be arrays which, conditional
on \(\mathcal H\), are independent of \(G\).  They may depend on
\(\mathcal H\) and may be dependent within and between arrays.
Suppose almost surely that
\begin{equation}
 \sum_i\E[Z_iZ_i^\top\mid\mathcal H]\preceq KI_r,
 \qquad \norm{Z_{i,a}}_U\le\delta,
 \qquad \norm{v_i}_U\le b\delta.
 \label{appx-an:conditional-budget}
\end{equation}
For a sufficiently large absolute constant \(A_b\),
\(h\ge A_b\delta\) implies
\begin{equation}
 \left|\sum_i\E D^{s+2}F(U+v_i)
         [Z_i,Z_i,Z_{i,1},\ldots,Z_{i,s}]\right|
 \le\frac{C_bK}{\sigma^2}\,
                 \delta^s h^{1-s}L^{3/2}.
 \label{appx-an:restoration-bound}
\end{equation}
\end{lemma}

\begin{proof}
Two Gaussian integrations by parts remove the repeated pair of
directions. We then expand the translated weights around the common
base \(U\), sum their matrix coefficients using the conditional
covariance budget, and apply Lemma~\ref{appx-an:lem:matrix}.
Condition on \(\mathcal H\) and on the \(i\)th array. The integrations
by parts give
\begin{align*}
 &\E_G D^{s+2}F(U+v_i)[Z_i,Z_i,Z_{i,1},\ldots,Z_{i,s}]\\
 &\qquad=\sigma^{-2}\E_G\left[
    \tr\{(GG^\top-I_r)Z_iZ_i^\top\}
    D^sF(U+v_i)[Z_{i,1},\ldots,Z_{i,s}]\right].
\end{align*}
For each partition \(\mathcal P=\{B_1,\ldots,B_k\}\) of
\(\{1,\ldots,s\}\), the derivative formula
\eqref{appx-eq:partition-derivative} contributes
\[
 h^{-s}(-1)^kk!\sum_{\alpha_1,\ldots,\alpha_k}
 F(x)\prod_{a=1}^k\pi_{\alpha_a}(x)
 \prod_{a=1}^k\prod_{q\in B_a}(u_{\alpha_a}\cdot Z_{i,q}).
\]
We now express every translated weight in terms of weights at \(U\).
For a shift \(v\), set
\[
 d_j(v)=e^{u_j\cdot v/h}-1,
 \qquad z_v(x)=\sum_j\pi_j(x)d_j(v),
 \qquad \eta=e^{b\delta/h}-1.
\]
Thus \(|d_j(v)|\le\eta\) and \(|z_v(x)|\le\eta\) when
\(\norm{v}_U\le b\delta\).  Directly from the selector definition,
\begin{equation}
 F(x+v)\prod_{a=1}^k\pi_{\alpha_a}(x+v)
 =F(x)\prod_{a=1}^k\pi_{\alpha_a}(x)
    e^{\sum_a u_{\alpha_a}\cdot v/h}
                    (1+z_v(x))^{-k-1}.
 \label{appx-an:weight-translation}
\end{equation}
For \(\eta<1\), use the absolutely convergent expansion
\[
 (1+z_v)^{-k-1}
 =\sum_{\ell\ge0}(-1)^\ell\binom{k+\ell}{\ell}z_v^\ell,
 \qquad
 z_v^\ell=\sum_{\beta_1,\ldots,\beta_\ell}
                \prod_{q=1}^\ell\pi_{\beta_q}(x)d_{\beta_q}(v).
\]
A term of order \(\ell\) has the common weight
\(F(U)\prod_a\pi_{\alpha_a}(U)\prod_q\pi_{\beta_q}(U)\).
Its summed matrix coefficient, conditional on \(\mathcal H\), is
\[
 B_{\alpha,\beta}=\sum_i\E\left[
  Z_iZ_i^\top
  \prod_{a=1}^k\prod_{q\in B_a}(u_{\alpha_a}\cdot Z_{i,q})
  e^{\sum_a u_{\alpha_a}\cdot v_i/h}
  \prod_{q=1}^\ell d_{\beta_q}(v_i)
       \,\middle|\,\mathcal H\right].
\]
For every deterministic vector \(z\),
\eqref{appx-an:conditional-budget} implies pointwise in \(\mathcal H\)
that
\[
 |z^\top B_{\alpha,\beta}z|
 \le\delta^se^{kb\delta/h}\eta^\ell
               \sum_i\E[(z\cdot Z_i)^2\mid\mathcal H]
 \le K\delta^se^{kb\delta/h}\eta^\ell\norm{z}^2.
\]
The matrix is symmetric, so its operator norm satisfies the same
bound.  This coefficient bound holds for every realization of \(\mathcal H\).
Apply trace duality,
\(|\tr(BM)|\le\norm{B}_{\mathrm{op}}\norm{M}_*\), then average over
\(\mathcal H\) and use \eqref{appx-an:matrix-bound} with \(k+\ell\)
weights. The concentration hypothesis is needed for this averaged matrix
estimate; no conditional concentration given \(\mathcal H\) is required.  If \(C_1\ge1\) is the constant in that matrix estimate, its
series factor is at most
\[
 e^{kb\delta/h}C_1^{k+1}
 \sum_{\ell\ge0}\binom{k+\ell}{\ell}(C_1\eta)^\ell
 =e^{kb\delta/h}C_1^{k+1}(1-C_1\eta)^{-k-1}
 \le C_b^{k+1},
\]
provided \(h\ge A_b\delta\).  Since \(s\le2\), the number of
partitions and their coefficients are bounded absolute constants.
Restoring \(\sigma^{-2}h^{-s}\) proves
\eqref{appx-an:restoration-bound}.

For completeness, the interchanges can first be justified with a
finite bound depending on \(r\).  The absolute face sums at order
\(\ell\) are at most \(\eta^\ell\), and their binomial series sums
to \((1-\eta)^{-k-1}\).  Moreover,
\(\sum_i\E\norm{Z_i}^2\le Kr<\infty\), so the Gaussian score is
integrable after summing over the arrays.  This finite trace is needed
only for justification and does not enter the final estimate.
\end{proof}

\subsection{The interpolation}
\label{appx-an:sec:path}

We can now define the comparison path.  Bernoulli thinning removes the
bounded rows while a Gaussian component supplies the lost covariance.
The choice of amplitude below also makes the derivative of the cubic
row coefficient vanish at the bounded endpoint; this will help control
the start of the path.

Use the independent support rows \(\zeta_i^0\) of
\eqref{appx-eq:support}, so
\[
 \E\zeta_i^0=0,\qquad
 \sum_i\E\zeta_i^0(\zeta_i^0)^\top=I_r,
 \qquad \norm{\zeta_i^0}_U\le\delta.
\]
For \(0\le\lambda\le1\), define
\begin{equation}
 \begin{gathered}
 a=1+\lambda,\qquad q=\frac{1-\lambda}{(1+\lambda)^2},
 \qquad B_{i,\lambda}\sim\operatorname{Bernoulli}(q),\\
 T_i=aB_{i,\lambda}\zeta_i^0,\qquad V_\lambda=\sum_iT_i,
 \qquad U_\lambda=V_\lambda+\sqrt\lambda G,
 \end{gathered}
 \label{appx-an:path}
\end{equation}
where the Bernoulli variables, the source rows, and \(G\) are
independent.  Only the law at each \(\lambda\) is used, so no
coupling of Bernoulli variables across \(\lambda\) is required.
The total covariance of \(U_\lambda\) is \(I_r\), since
\(qa^2=1-\lambda\).  For a rectangle \(A\), fix its selector
\(F=F_{A,h}\) on the support and put
\[
 P_A(\lambda)=\E F(U_\lambda),\qquad H_i=U_\lambda-T_i.
\]
Let \(\zeta_i,\widetilde\zeta_i\) be two independent fresh copies of
\(\zeta_i^0\), independent of the actual path array and \(G\).

\begin{lemma}[Exact fourth-order identity]\label{appx-an:lem:identity}
For \(0<\lambda<1\),
\begin{equation}
 P_A'(\lambda)
 =-\frac\lambda3\sum_i\E D^3F(U_\lambda)
                    [\zeta_i,\zeta_i,\zeta_i]
       +\sum_i(\mathcal A_i+\mathcal B_i+\mathcal C_i),
 \label{appx-an:fourth-identity}
\end{equation}
where
\begin{align}
 \mathcal A_i
 &=\int_0^1 k_\lambda(s)
       \E D^4F(H_i+sa\zeta_i)
                 [\zeta_i,\zeta_i,\zeta_i,\zeta_i]\dd s,
 \label{appx-an:A}\\
 k_\lambda(s)
 &=\frac a6(1-s)^2\{(3-\lambda)s-2\lambda\},
 \label{appx-an:kernel}\\
 \mathcal B_i
 &=\frac{1-\lambda}{2}\int_0^1(1-s)
       \E D^4F(H_i+sa\widetilde\zeta_i)
          [\zeta_i,\zeta_i,\widetilde\zeta_i,\widetilde\zeta_i]
       \dd s,
 \label{appx-an:B}\\
 \mathcal C_i
 &=\frac{\lambda(1-\lambda)}{3a}\int_0^1
       \E D^4F(H_i+sa\widetilde\zeta_i)
          [\zeta_i,\zeta_i,\zeta_i,\widetilde\zeta_i]\dd s.
 \label{appx-an:C}
\end{align}
The three displayed fourth-order terms account for the entire remainder.
\end{lemma}

\begin{proof}
Differentiate the Bernoulli row operators and combine their quadratic
terms with the Gaussian heat derivative. Then move the cubic terms
to the full base; the three resulting fourth-order terms are the ones
in the statement. Set \(M_i=\E\zeta_i\zeta_i^\top\) and
\(K_i^{(3)}=\E\zeta_i^{\otimes3}\).  For a smooth test function
\(f\), the row operator is
\[
 Q_if(x)=(1-q)f(x)+q\E f(x+a\zeta_i).
\]
Differentiation in \(\lambda\) gives
\[
 \dot Q_if(x)=q'\{\E f(x+a\zeta_i)-f(x)\}
                  +q\E Df(x+a\zeta_i)[\zeta_i].
\]
Apply Taylor's formula through degree three to the difference, and
through degree two to the derivative.  Centering removes the linear
terms.  Since \((qa^2)'=-1\) and \((qa^3)'=-2\lambda\), the result is
\[
 \dot Q_if
 =-\tfrac12 M_i:D^2f
   -\tfrac\lambda3 K_i^{(3)}:D^3f
   +\int_0^1k_\lambda(s)
        \E D^4f(\,\cdot+sa\zeta_i)[\zeta_i^4]\dd s.
\]
Here the notation \([\zeta_i^4]\) means four repeated arguments,
and the fourth-order kernel is exactly
\[
 \frac{q'a^4}{6}(1-s)^3+\frac{qa^3}{2}(1-s)^2
 =\frac a6(1-s)^2\{(3-\lambda)s-2\lambda\}.
\]
The Gaussian heat derivative contributes
\(\tfrac12\sum_iM_i:\E D^2F(U_\lambda)\).  Combining this with
the quadratic row terms leaves
\[
 \frac12\sum_i M_i:
             \E\{D^2F(U_\lambda)-D^2F(H_i)\}.
\]
Conditional on \(H_i\), the deleted row is a centered independent
Bernoulli row.  Its second-order Taylor formula therefore identifies
this expression with \(\sum_i\mathcal B_i\), using \(qa^2=1-\lambda\).
Finally,
\[
 -\frac\lambda3 K_i^{(3)}:\E D^3F(H_i)
 =-\frac\lambda3 K_i^{(3)}:\E D^3F(U_\lambda)
   +\frac\lambda3 K_i^{(3)}:
          \E\{D^3F(U_\lambda)-D^3F(H_i)\}.
\]
A first-order integral expansion of the last difference is exactly
\(\mathcal C_i\), because \(qa=(1-\lambda)/a\).  The remaining
fourth-order row term is \(\mathcal A_i\).  This proves the identity.
All differentiations are legitimate for fixed \(h>0\): the number
of rows is finite, the rows are bounded in the covariance support,
and the selector and its derivatives are bounded.
\end{proof}

The identity gives two useful estimates for the path derivative.
Near \(\lambda=0\), the boundary concentration estimate controls all
derivatives directly.  Once the Gaussian component is large enough,
integration by parts gives a sharper bound.

\begin{proposition}[Two bounds for the path derivative]
\label{appx-an:prop:derivative}
There are sufficiently large absolute constants \(A,d_0\) such that,
for \(h\ge A\delta\), uniformly over rectangles \(A\),
\begin{equation}
 |P_A'(\lambda)|
 \le C\sqrt L\left(\frac{\lambda\delta}{h^2}
                         +\frac{\delta^2}{h^3}\right),
 \qquad 0<\lambda<1.
 \label{appx-an:early}
\end{equation}
For \(d_0\delta^2L\le\lambda<1\), one also has
\begin{equation}
 |P_A'(\lambda)|
 \le CL^{3/2}\left(\delta+\frac{\delta^2}{\lambda h}\right).
 \label{appx-an:late}
\end{equation}
More specifically, for the fourth-order remainder
\(R_{4,A,h}(\lambda)=\sum_i(\mathcal A_i+\mathcal B_i+\mathcal C_i)\),
\begin{equation}
 |R_{4,A,h}(\lambda)|\le\frac{C\delta^2}{\lambda h}L^{3/2},
 \qquad
 \left|\sum_i\mathcal C_i\right|
       \le\frac{C\delta^2}{h}L^{3/2}.
 \label{appx-an:remainder}
\end{equation}
\end{proposition}

\begin{proof}
Both estimates begin with the same concentration bound for the full path:
\begin{equation}
 \E F(U_\lambda)\theta(U_\lambda)\le Ch\sqrt L.
 \label{appx-an:path-weight}
\end{equation}
For \(\lambda\le1/2\), its bounded independent part has every face
variance \(1-\lambda\ge1/2\) and row envelope at most \(2\delta\).
For \(\lambda\ge1/2\), its independent Gaussian part has every
face variance at least \(1/2\).  Lemmas~\ref{appx-lem:concentration}
and \ref{appx-lem:envelope} give \eqref{appx-an:path-weight} in both cases.

To prove \eqref{appx-an:early}, work temporarily in the original
coordinates.  Put \(\xi_i=B\zeta_i\) and
\(\widetilde\xi_i=B\widetilde\zeta_i\).  These fresh directions
are independent of the full path \(BU_\lambda\), have maximum norm
at most \(\delta\), and satisfy
\(\sum_i\E\xi_{ij}^2=1\) for every coordinate \(j\).  For any
coordinate triple,
\[
 \sum_i\E|\xi_{ij}\xi_{ik}\xi_{i\ell}|
 \le\delta\sum_i(\E\xi_{ij}^2\E\xi_{ik}^2)^{1/2}
 \le\delta.
\]
Likewise, every fourth-order coefficient appearing in
\eqref{appx-an:A}, \eqref{appx-an:B}, or \eqref{appx-an:C}, after taking absolute
values and summing over \(i\), is at most \(\delta^2\).  Retain
the repeated fresh pair, bound the other two factors by \(\delta^2\),
and apply the same Cauchy--Schwarz argument to that pair.

Each deleted-row base is a translate of the full path:
\[
 H_i+sa\zeta_i=U_\lambda+v_i,
 \qquad v_i=sa\zeta_i-T_i,
\]
and similarly with \(\widetilde\zeta_i\).  Its shift has
\(U\)-norm at most \(4\delta\).  Its original-coordinate image is
\(Bv_i=sa\xi_i-BT_i\), which directly has maximum norm at most
\(4\delta\); the same holds for the other fresh copy.  Take the
supremum over all such shifts in the
local derivative envelope \eqref{appx-eq:local-envelope}.  The moment
coefficients above can then be separated from the full base, since
the fresh directions are independent of that base.  The cubic term
in \eqref{appx-an:fourth-identity} is bounded by
\(C\lambda\delta h^{-3}\E F\theta\); the three fourth-order terms
together are bounded by \(C\delta^2h^{-4}\E F\theta\).
The absolute integrals of their scalar kernels on \([0,1]\) are
bounded uniformly in \(\lambda\).  Substituting
\eqref{appx-an:path-weight} proves \eqref{appx-an:early}.

For \eqref{appx-an:late}, we construct a base to which
Lemma~\ref{appx-an:lem:restoration} applies. That lemma needs both a
concentrated base and an independent Gaussian component for integration
by parts. Near \(\lambda=1\), the bounded part no longer supplies the
first property, so we split the Gaussian component.  If \(\lambda\le1/2\), take
\(V_0=V_\lambda\), \(\sigma^2=\lambda\), and let \(\mathcal H\)
contain the entire actual path array.  If \(\lambda\ge1/2\), split
\(\sqrt\lambda G\) into two independent Gaussian vectors with
covariance \((\lambda/2)I_r\); put one of them into \(V_0\) and
\(\mathcal H\), and use the other as \(\sigma G\), so that
\(\sigma^2=\lambda/2\).  The first base satisfies
\eqref{appx-an:concentration} by its bounded component of face variance
at least \(1/2\).  The second satisfies it by its Gaussian component
of face variance at least \(1/4\).  Thus in either case
\(\sigma^2\asymp\lambda\), and
\(\lambda\ge d_0\delta^2L\) ensures
\(\sigma\ge\delta\sqrt L\).

In the cubic term use Lemma~\ref{appx-an:lem:restoration} with
\(s=1\), \(Z_i=Z_{i,1}=\zeta_i\), and \(v_i=0\).
Freshness gives the conditional covariance budget
\(\sum_i\E[\zeta_i\zeta_i^\top\mid\mathcal H]=I_r\).
Consequently,
\begin{equation}
 \left|\frac\lambda3\sum_i\E D^3F(U_\lambda)
                    [\zeta_i,\zeta_i,\zeta_i]\right|
 \le C\delta L^{3/2}.
 \label{appx-an:cubic-bound}
\end{equation}
For each fourth-order term, restore the actual row as above and
apply the same lemma with \(s=2\), \(Z_i=\zeta_i\), and \(b=4\).
The remaining two directions are, respectively,
\((\zeta_i,\zeta_i)\),
\((\widetilde\zeta_i,\widetilde\zeta_i)\), and
\((\zeta_i,\widetilde\zeta_i)\).  Each has \(U\)-norm at most
\(\delta\), and the same conditional covariance budget applies.
Every summed derivative integrand is therefore bounded in absolute
value by \(C\delta^2L^{3/2}/(\lambda h)\).
The scalar kernels for \(\mathcal A_i,\mathcal B_i\) have bounded
absolute integrals, and the coefficient of \(\mathcal C_i\) has
an additional factor \(\lambda\).  This proves
\eqref{appx-an:remainder}; together with \eqref{appx-an:cubic-bound}
it proves \eqref{appx-an:late}.
\end{proof}

Finally, the path is continuous in law at the endpoints:
\(BU_0\) has the law of \(W\), and \(U_1\) is standard Gaussian
on the support.  Since \(F\) is bounded and continuous,
\(P_A(0)=\E F_{A,h}(W)\) and
\(P_A(1)=\E F_{A,h}(Z)\), with the selector interpreted in the
appropriate coordinates.  The bound \eqref{appx-an:early} is integrable
on the whole interval and justifies integration of the derivative
identity to both endpoints.


% Source: concentration_proof
\section{Verification of the concentration input}\label{appx-app:concentration}

We prove Lemma~\ref{appx-lem:concentration} by resampling a small proportion
of the rows. Away from a tail event, a sum near a rectangle face has a
fixed positive chance of crossing it; noise sensitivity bounds the total probability of a
change in membership. We follow~\cite[Lemma 2.1]{AL}. The external
input is Kane's noise-sensitivity bound~\cite[Corollary 1]{Kane}: if \(f\) is the indicator of an intersection
of \(k\) affine halfspaces on a uniform Boolean cube, and
\(\varepsilon^\theta\) is obtained by independently flipping signs
with probability \(\theta\in(0,1/2]\), then
\begin{equation}
 \Pp\{f(\varepsilon)\ne f(\varepsilon^\theta)\}
 \le C\sqrt{\theta\log(ek)}.
 \label{appx-eq:kane}
\end{equation}

First normalize to \(v=1\) and omit the independent shift. Thus
\(S=\sum_i\zeta_i\) has independent centered rows,
\(\norm{\zeta_i}_\infty\le\eta\), and \(\Var(S_j)=1\) for every \(j\).
Independently replace each row by a fresh copy with probability
\(\theta\), obtaining \(S^\theta\).
To apply \eqref{appx-eq:kane}, generate two independent candidate copies
\(\zeta_i^+,\zeta_i^-\), and choose between them using a uniform sign.
Conditional on all candidate pairs, membership in a rectangle is an
intersection of at most \(2p\) affine halfspaces in these signs.
A sign flip selects the other independent candidate, so averaging
\eqref{appx-eq:kane} gives
\begin{equation}
 \Pp\{\ind_Q(S)\ne\ind_Q(S^\theta)\}
 \le C\sqrt{\theta L}
 \label{appx-eq:resampling}
\end{equation}
for every rectangle \(Q\). Open or closed faces cause no difficulty on
a finite Boolean cube, since thresholds can be perturbed without changing
the Boolean function.

To show that points near a face can leave the rectangle, condition on
the original rows. Use the equivalent representation
\(S^\theta=\sum_i[(1-B_i)\zeta_i+B_i\zeta_i']\), where the
\(B_i\sim\operatorname{Bernoulli}(\theta)\) and the fresh rows are
independent of \(\mathcal F=\sigma(\zeta_1,\ldots,\zeta_n)\).
For \(D=S^\theta-S\), each coordinate obeys
\[
 D_j+\theta S_j=\sum_iY_{ij},\qquad
 Y_{ij}=B_i(\zeta_{ij}'-\zeta_{ij})+\theta\zeta_{ij}.
\]
The \(Y_{ij}\) are conditionally independent and centered,
\(|Y_{ij}|\le2\eta\), and
\[
 v_j:=\sum_i\E(Y_{ij}^2\mid\mathcal F)
 =\theta+\theta(1-\theta)\sum_i\zeta_{ij}^2\ge\theta.
\]
Also, \(\sum_i\E(|Y_{ij}|^3\mid\mathcal F)\le2\eta v_j\).
The scalar Berry--Esseen inequality therefore yields, almost surely,
\begin{equation}
 \sup_x\left|
 \Pp\left\{\frac{D_j+\theta S_j}{\sqrt{v_j}}\le x
 \,\middle|\,\mathcal F\right\}-\Phi(x)\right|
 \le C\eta/\sqrt\theta.
 \label{appx-eq:conditional-be}
\end{equation}

Choose the resampling variance large enough to cross a layer of width
\(d\), while keeping the Berry--Esseen error small. Let \(\partial_d Q\)
be the points in \(Q\) within distance \(d\) of a finite face. Put
\(q=d+\eta\) and \(\theta=Kq^2\), with \(K\) a large absolute
constant. We may assume \(q\sqrt L\) is sufficiently small, since
otherwise the desired estimate follows from the bound by one. This
assumption also ensures \(\theta\le1/2\).
On
\(\{S\in\partial_d Q,\ \norm S_\infty\le1/(4\sqrt\theta)\}\),
select the first face within distance \(d\). This choice is
\(\mathcal F\)-measurable, so \eqref{appx-eq:conditional-be} applies to
its coordinate and outward direction. Conditional on \(\mathcal F\), the standardized displacement needed to cross that
face outward is at most
\[
 d/\sqrt\theta+\sqrt\theta|S_j|
 \le K^{-1/2}+1/4.
\]
The error in \eqref{appx-eq:conditional-be} is at most \(C K^{-1/2}\).
Choose \(K\) large enough that the conditional outward-crossing
probability is bounded below by an absolute positive constant.
Crossing forces \(S^\theta\notin Q\). Equation~\eqref{appx-eq:resampling}
then gives
\[
 \Pp(S\in\partial_d Q)
 \le Cq\sqrt L
 +\Pp\{\norm S_\infty>1/(4\sqrt\theta)\}.
\]
Bernstein's inequality and a union bound imply
\[
 \Pp\{\norm S_\infty>1/(4\sqrt\theta)\}
 \le2\exp(L-c_1/\theta).
\]
Here \(\eta/\sqrt\theta\le K^{-1/2}\) controls the bounded-summand
term in Bernstein's denominator. Decreasing the smallness threshold for
\(q\sqrt L\) ensures \(L\le c_1/(2\theta)\), and therefore this
tail is at most \(Cq\sqrt L\). We have proved
\begin{equation}
 \Pp(S\in\partial_d Q)\le C(d+\eta)\sqrt L.
 \label{appx-eq:inner-layer}
\end{equation}

Every point of \(A^a\setminus A^{-a}\) lies within distance \(2a\)
of a finite face of \(A^a\). If an interval of \(A^{-a}\) is empty,
the corresponding interval of \(A^a\) has length at most \(4a\), so
the same assertion holds. Apply \eqref{appx-eq:inner-layer} with
\(Q=A^a\) and \(d=2a\). An independent shift is handled by
conditioning, and equal coordinate variances \(v\) by dividing by
\(\sqrt v\). This proves the bounded-sum part of
Lemma~\ref{appx-lem:concentration}.

For the Gaussian assertion, factor its covariance as
\(\Sigma=\sum_{k=1}^r v_kv_k^\top\), and approximate the Gaussian by
\[
 G_m=\frac1{\sqrt m}\sum_{k=1}^r\sum_{l=1}^m v_k\varepsilon_{kl},
\]
where the \(\varepsilon_{kl}\) are independent uniform signs.
The covariance remains \(\Sigma\), while the maximum summand norm
tends to zero. Every coordinate has variance \(v>0\), so each finite coordinate face
has Gaussian probability zero, even if \(\Sigma\) is singular. Each
fixed rectangle boundary layer is therefore a continuity set for the
Gaussian law. The finite-dimensional CLT passes the bounded-sum estimate
to the limit. Conditioning again permits independent shifts.



% Source: complete_integration_proofs
\section{Integrating the comparison}\label{sec:main-proof}
We integrate the early derivative bound until the Gaussian component
makes the later bound more effective. The two fourth-order estimates
agree at \(\lambda=Lh^2\); this switching time produces the logarithm
in the smoothed comparison. We then balance that comparison against
the endpoint smoothing error.

\begin{proposition}[Smoothed comparison at every admissible scale]
\label{prop:moving-cutoff}
There are absolute constants \(A,C\) such that under \eqref{eq:setup},
for every \(h\ge A\delta\),
\begin{equation}
 M_h\le CL^{3/2}\left[
 \delta+\frac{\delta^2}{h}\log_+\frac1{Lh^2}\right].
 \tag{S*}\label{eq:smooth-moving}
\end{equation}
The estimate is uniform over all finite threshold systems formed from
signed coordinate faces, including inconsistent systems.
\end{proposition}

\begin{proof}
Increase \(A\) so that \eqref{an:early} and \eqref{an:late} apply and
\(A^2\ge d_0\). Put \(\tau_h=Lh^2\), the time at which their
fourth-order terms agree. Since \(h\ge A\delta\), we have
\(\tau_h\ge d_0\delta^2L\). If \(\tau_h<1\), the later estimate
is therefore available throughout \([\tau_h,1]\). Integrating the
early estimate gives
\begin{align*}
 |P_A(\tau_h)-P_A(0)|
 &\le C\sqrt L\left(
 \frac{\delta\tau_h^2}{h^2}
 +\frac{\delta^2\tau_h}{h^3}\right)\\
 &=C\delta L^{3/2}
 \left(Lh^2+\frac\delta h\right)
 \le C\delta L^{3/2}.
\end{align*}
The later estimate gives
\[
 |P_A(1)-P_A(\tau_h)|
 \le CL^{3/2}\left[
 \delta+\frac{\delta^2}{h}\log\frac1{\tau_h}\right].
\]
The constant cubic contribution integrates to \(C\delta L^{3/2}\);
the fourth-order factor \(1/\lambda\) gives the logarithm. Combining
the two time ranges and taking the supremum over thresholds proves
\eqref{eq:smooth-moving} when \(\tau_h<1\).

If \(Lh^2\ge1\), integrate the early estimate over the whole path:
\[
 |P_A(1)-P_A(0)|
 \le C\sqrt L\left(\frac\delta{h^2}
 +\frac{\delta^2}{h^3}\right)
 =C\delta L^{3/2}\frac{1+\delta/h}{Lh^2}
 \le C\delta L^{3/2}.
\]
The logarithmic term in \eqref{eq:smooth-moving} is zero in this case.
Continuity at the interpolation endpoints permits integration up to
\(0\) and \(1\). All bounds are uniform in the thresholds.
\end{proof}

\begin{proof}[Proof of Theorem~\ref{thm:optimized}]
Combining Proposition~\ref{prop:moving-cutoff} with
\eqref{eq:endpoint-min} yields, for every \(h\ge A\delta\),
\begin{equation}
 \rh(W,Z)\le C\left[
 \delta L^{3/2}
 +\frac{\delta^2}{h}L^{3/2}\log_+\frac1{Lh^2}
 +h\sqrt L\,\kappa\right].
 \label{eq:variational}
\end{equation}
The first term is independent of \(h\). To balance the fourth-order
comparison error against the endpoint error, choose
\begin{equation}
 h=A\delta\sqrt{\frac{LE_\kappa}{\kappa}}.
 \label{eq:optimal-scale}
\end{equation}
Because \(\kappa\le L\) and \(E_\kappa\ge1\), this scale is at least
\(A\delta\). Moreover,
\[
 \log_+\frac1{Lh^2}
 =2\log_+\frac{\sqrt\kappa}{A\delta L\sqrt{E_\kappa}}
 \le2\log_+\frac{\sqrt\kappa}{\delta L}
 \le2E_\kappa.
\]
The endpoint term in \eqref{eq:variational} equals
\(A\delta L\sqrt{\kappa E_\kappa}\), while its logarithmic term is at
most \((2/A)\delta L\sqrt{\kappa E_\kappa}\).
This proves \eqref{eq:optimized} after taking the minimum with one.
The version with \(D_L\) follows from \(E_\kappa\le D_L\).
\end{proof}


\begin{proof}[Proof of Theorem~\ref{thm:unrestricted}]
Since \(\kappa\le L\), \(E_\kappa\le D_L\), and \(D_L\ge1\),
\[
 L^{3/2}+L\sqrt{\kappa E_\kappa}
 \le 2L^{3/2}\sqrt{D_L}.
\]
Apply Theorem~\ref{thm:optimized}.
\end{proof}


% Source: zero_skewness_proof
\subsection{When the cubic term cancels}\label{zs:section}
\begin{proposition}[Zero-skewness smoothed comparison]
\label{zs:smooth-proposition}
Under \eqref{eq:setup} and \eqref{eq:zero-third}, there is an absolute
constant \(A\) such that, for every \(h\ge A\delta\),
\begin{equation}
 M_h\le \frac{C\delta^2L^{3/2}}h
 \left[1+\log_+\frac1{Lh^2}\right].
 \label{zs:smooth}
\end{equation}
The estimate is uniform over the signed coordinate threshold systems
used in the smoothing endpoint lemmas.
\end{proposition}

\begin{proof}
Use \eqref{zs:identity} to cancel the common-base cubic term before
estimating the derivative. The separate fourth-order bounds
\eqref{eq:remainder-early} and \eqref{an:remainder} then give
\begin{align}
 |P_A'(\lambda)|
 &\le \frac{C\delta^2\sqrt L}{h^3},
 &&0<\lambda<1,\label{zs:early}\\
 |P_A'(\lambda)|
 &\le \frac{C\delta^2L^{3/2}}{\lambda h},
 &&d_0\delta^2L\le\lambda<1.\label{zs:late}
\end{align}
Increase \(A\) so that \(A^2\ge d_0\), and put \(\tau_h=Lh^2\).
If \(\tau_h<1\), then \(\tau_h\ge d_0\delta^2L\). Integrating
\eqref{zs:early} up to \(\tau_h\), and \eqref{zs:late} after it,
yields
\begin{align*}
 |P_A(1)-P_A(0)|
 &\le \frac{C\delta^2\sqrt L}{h^3}\tau_h
       +\frac{C\delta^2L^{3/2}}h\log\frac1{\tau_h}\\
 &=\frac{C\delta^2L^{3/2}}h
       \left[1+\log\frac1{Lh^2}\right].
\end{align*}
If \(Lh^2\ge1\), integrate \eqref{zs:early} on the whole path and
use \(\sqrt L/h^3\le L^{3/2}/h\). Endpoint continuity, already
established after Proposition~\ref{an:prop:derivative}, justifies
integration to \(0\) and \(1\). Taking the supremum over thresholds
proves \eqref{zs:smooth}.
\end{proof}

\begin{proof}[Proof of Theorem~\ref{thm:zero-skew}]
Combine Proposition~\ref{zs:smooth-proposition} with the endpoint
bound \eqref{eq:endpoint-min} to obtain, for every \(h\ge A\delta\),
\begin{equation}
 \rh(W,Z)\le C\left\{
 \frac{\delta^2L^{3/2}}h
 \left[1+\log_+\frac1{Lh^2}\right]+h\sqrt L\,\kappa\right\}.
 \label{zs:variational}
\end{equation}
To balance the comparison error against the endpoint error, choose
\eqref{eq:optimal-scale} again:
\[
 h=A\delta\sqrt{\frac{LE_\kappa}{\kappa}}.
\]
It is admissible because \(\kappa\le L\) and \(E_\kappa\ge1\).
As before,
\[
 \log_+\frac1{Lh^2}
 =2\log_+\frac{\sqrt\kappa}{A\delta L\sqrt{E_\kappa}}
 \le2E_\kappa.
\]
The endpoint term in \eqref{zs:variational} equals
\(A\delta L\sqrt{\kappa E_\kappa}\); its comparison term is at most
\[
 \frac{\delta L\sqrt\kappa}{A\sqrt{E_\kappa}}(1+2E_\kappa)
 \le\frac3A\delta L\sqrt{\kappa E_\kappa}.
\]
Taking the minimum with one proves \eqref{eq:zero-skew}.

\end{proof}


% Source: subexponential_extension_proof
\section{Sub-exponential coordinates}\label{app:subexponential}
\begin{proof}[Proof of Corollary~\ref{cor:subexponential}]
We clip each coordinate, apply the bounded theorem after marginal
standardization, and transfer the result back to \(W_n\). This requires
control of both the discarded sum and the Gaussian covariance change.
Write \(\ell=\log(en)\) and \(b=K\ell/\sqrt n\). The unrestricted
bound follows from the bound by one, after adjusting its absolute
constant, when \(b\ge1/16\), so assume \(b<1/16\).

\emph{Step 1: Clip the rows and bound the residual sum.}
Put \(u=4\ell\), \(T=Ku\), and define coordinatewise
\[
 C_T(x)=\operatorname{sgn}(x)\min\{|x|,T\},\qquad
 Y_i=C_T(X_i)-\E C_T(X_i),\qquad R_i=X_i-C_T(X_i).
\]
Here \(Y_i\) is the centered clipped row and \(R_i\) is the discarded
part. Their normalized sums satisfy
\[
 W_T=n^{-1/2}\sum_iY_i,\qquad
 V=n^{-1/2}\sum_i(R_i-\E R_i),\qquad W_n=W_T+V.
\]
The exponential-moment assumption gives, for \(t\ge0\),
\[
 \Pp(|R_{ij}|>t)\le2e^{-u-t/K},\qquad
 \E|R_{ij}|^k\le2k!K^ke^{-u}\quad(k\ge1).
\]
Consequently,
\begin{equation}
 \max_j\Var(V_j)\le4K^2e^{-u},\qquad
 \E|R_{ij}-\E R_{ij}|^k
 \le\frac{k!}{2}(16K^2e^{-u})(2K)^{k-2}
 \quad(k\ge2).
 \label{tail:residual-moments}
\end{equation}
For the second estimate use
\(\E|R-\E R|^k\le2^k\E|R|^k\).
Thus Bernstein's inequality applies to each coordinate of \(V\) with
variance parameter \(16K^2e^{-u}\) and scale \(2K/\sqrt n\).
With \(s=\log(4np)\), a coordinate union bound yields
\begin{equation}
 \Pp(\|V\|_\infty>r_1)\le\frac1{2n},\qquad
 r_1=\sqrt{32}\,Ke^{-u/2}\sqrt s+\frac{2Ks}{\sqrt n}.
 \label{tail:sum-coupling}
\end{equation}
We control the discarded sum rather than requiring all \(np\) entries
to lie below \(T\). This is why the cutoff needs only \(\log n\).

\emph{Step 2: Apply the bounded comparison.}
Set \(d_j=\sqrt{\Var(W_{T,j})}\) and \(\tau=2Ke^{-u/2}\).
Since each coordinate of \(W_n\) has variance one, the triangle
inequality in \(L^2\) gives
\[
 |d_j-1|\le\tau,
 \qquad \frac\tau b=\frac{2\sqrt n\,e^{-2\ell}}\ell
 \le\frac2{e^2}.
\]
Hence \(d_j\ge1/2\). Since \(\|Y_i\|_\infty\le2T\), the
independent centered rows of the marginally standardized \(W_T\)
have envelope at most
\[
 \frac{2T}{\sqrt n\min_jd_j}\le16b.
\]
Let \(Z_T\) be Gaussian with covariance \(\Cov(W_T)\).
Positive diagonal scaling preserves rectangle distance. Applying
Theorem~\ref{thm:unrestricted} with envelope \(16b\) gives
\begin{equation}
 \rh(W_T,Z_T)\le C\min\left\{1,
 bL^{3/2}\sqrt{1+\log_+\frac1{b\sqrt L}}\right\}.
 \label{tail:bounded-comparison}
\end{equation}
The factor 16 is absorbed into the absolute constant, since increasing
the envelope decreases the logarithmic factor.

\emph{Step 3: Restore the covariance.}
Construct a jointly centered Gaussian vector \((Z_T,H)\) with the
block covariance of \((W_T,V)\). This covariance is positive
semidefinite, even when singular. Set \(Z_n=Z_T+H\) on this coupling.
Then \(Z_n\) has exactly the covariance of \(W_n\), while
\(\max_j\Var(H_j)\le\tau^2\). Gaussian tails therefore give
\begin{equation}
 \Pp(\|H\|_\infty>r_2)\le\frac1{2n},\qquad
 r_2=\sqrt8\,Ke^{-u/2}\sqrt s.
 \label{tail:Gaussian-coupling}
\end{equation}
The Gaussian components need not be independent.

For any rectangle \(A\), the two coupling bounds give
\begin{align*}
 \Pp(W_n\in A)
 &\le\Pp(W_T\in A^{r_1})+\frac1{2n}\\
 &\le\Pp(Z_T\in A^{r_1})+\rh(W_T,Z_T)+\frac1{2n}\\
 &\le\Pp(Z_n\in A^{r_1+r_2})+\rh(W_T,Z_T)+\frac1n.
\end{align*}
Inward expansions give the reverse bound
\(\Pp(W_n\in A)\ge\Pp(Z_n\in A^{-(r_1+r_2)})
-\rh(W_T,Z_T)-1/n\). Applying Gaussian rectangle concentration at
the unit-marginal target \(Z_n\), from Lemma~\ref{lem:concentration}, yields
\begin{equation}
 \rh(W_n,Z_n)\le\rh(W_T,Z_T)
       +C(r_1+r_2)\sqrt L+\frac1n.
 \label{tail:transfer}
\end{equation}
No independence between the two vectors within either coupling is
required for this sandwich argument.
The transfer error has order \(bL^{3/2}\): since
\(s\le2\ell L\) and \(K\ge1\),
\[
 (r_1+r_2)\sqrt L+\frac1n
 \le CK\left[e^{-2\ell}\sqrt{Ls}
                     +\frac{s\sqrt L}{\sqrt n}\right]+\frac1n
 \le CbL^{3/2}.
\]
Combine this with \eqref{tail:bounded-comparison} and the bound by one
to prove \eqref{eq:subexponential}.

\emph{Step 4: Preserve a fixed correlation gap.}
For \eqref{eq:subexponential-gap}, it suffices to consider
\(b<\min\{1/16,c_0/4\}\); otherwise its right-hand side is at
least one after increasing \(C_{c_0}\).
Under the Gaussian coupling, both \(Z_n\) and the standardized clipped
vector \((Z_{T,j}/d_j)_{j\le p}\) have unit coordinate variances, and
\[
 \|Z_{T,j}/d_j-Z_{n,j}\|_{L^2}\le2\tau.
\]
Writing the difference of two covariance entries as a sum of two
products and applying Cauchy--Schwarz shows that the covariance matrices
differ entrywise by at most \(4\tau\). Since
\(4\tau\le(8/e^2)b<c_0/2\), the standardized clipped covariance
has correlation gap at least \(c_0/2\).
Theorem~\ref{thm:optimized} applied with envelope \(16b\) now gives
\[
 \rh(W_T,Z_T)\le C_{c_0}\min\left\{1,
 b\left[L^{3/2}+L\sqrt{1+\log_+\frac1{bL}}\right]\right\}.
\]
The transfer error in \eqref{tail:transfer} is bounded by the
\(CbL^{3/2}\) term. This proves \eqref{eq:subexponential-gap}.
\end{proof}

Coordinatewise odd clipping preserves central symmetry of each row,
but need not preserve cancellation of third tensors across rows.
Moreover the transfer cost above contains \(bL^{3/2}\), so the
zero-skewness refinement cannot simply be read off by deleting the
cubic term from the displayed sub-exponential bound.

\paragraph{Comparison with the locally nondegenerate result.}
The companion paper~\cite[Corollary~1]{Local} assumes the same
coordinatewise \(\psi_1\) condition, an averaged fourth-moment bound
\(\max_jn^{-1}\sum_i\E|X_{ij}|^4\le B_4^2\), and uniformly positive
eigenvalues for principal covariance submatrices of order at most four.
When these principal-submatrix eigenvalues are bounded below by \(a>0\), its bound is
\[
 \frac{C_a}{\sqrt n}\left[B_4L^{3/2}
                   +K\log(enp)\sqrt L\right],
\]
capped at one. For fixed \(K,B_4,a\) and \(p\asymp n^q\), \(q>0\),
this is \(O(n^{-1/2}(\log n)^{3/2})\), whereas
\eqref{eq:subexponential} is \(O(n^{-1/2}(\log n)^3)\).
The latter requires neither a fourth-moment hypothesis in addition to
the tail bound nor a local covariance assumption. If \(K\) varies,
the elementary fourth-moment consequence is \(B_4\le C K^2\);
the tail bound alone does not justify taking \(B_4=K\).


\end{document}
